\documentclass[12pt,a4paper]{scrreprt}

\usepackage{iftex}

% --- Encoding + PDF Unicode strategy ---
% Goal: (1) correct rendering in the PDF, (2) clean copy&paste/search, (3) clean bookmarks/metadata.
% - pdfLaTeX: uses inputenc/fontenc + ToUnicode mappings.
% - LuaLaTeX/XeLaTeX: uses native Unicode (fontspec).
\ifPDFTeX
  % Improves copy&paste/search by providing ToUnicode mappings (fixes e.g. ff/fi/fl ligatures
  % turning into control characters in extracted text).
  \usepackage{cmap}
  \usepackage[utf8]{inputenc}
  \usepackage[T1]{fontenc}
  \usepackage{lmodern}
  \usepackage{microtype}

  % Generate Unicode mappings for glyphs in the PDF (pdfTeX).
  \input{glyphtounicode}
  \pdfgentounicode=1
\else
  \usepackage{fontspec}
  \defaultfontfeatures{Ligatures=TeX}
  \usepackage{microtype}
\fi

\usepackage[ngerman]{babel}

\usepackage{graphicx}
\usepackage{tikz}
\usepackage{amsmath}
\usepackage{siunitx}
\usepackage{csquotes}

\sisetup{
  locale = DE,
  per-mode = symbol
}

\usepackage[
  backend=biber,
  style=authoryear,
  sorting=nyt,
  maxcitenames=2,
  maxbibnames=99,
  doi=true,
  url=true
]{biblatex}

\addbibresource{references.bib}

\usepackage{amsthm}
\newtheorem{hypothesis}{Hypothese}

% Hyperlinks, PDF metadata, bookmarks.
\usepackage[unicode]{hyperref}
\ifPDFTeX
  % Better Unicode coverage for PDF strings (bookmarks/metadata) under pdfLaTeX.
  % Optional: not all TeX installations include it.
  \IfFileExists{psdextra.sty}{\usepackage{psdextra}}{}
\fi
\usepackage{bookmark}

\csname title\endcsname{Einfluss der Entwicklerkompetenz auf die Effektivität von KI-Assistenzsystemen:\\
Eine empirische Untersuchung von Vibe-Coding in der Unternehmenssoftwareentwicklung}
\author{Cihad Gök}
\date{Dezember 2025}

\begin{document}

\maketitle

\tableofcontents
\clearpage

\chapter{Einleitung}
\label{chap:introduction}

\section{Motivation}
Große Sprachmodelle (Large Language Models, LLMs) werden zunehmend in der Softwareentwicklung eingesetzt und unterstützen Entwurf, Implementierung und Wartung.
Im Folgenden wird ein solches \emph{interaktives} Assistenzsystem als \emph{KI-Assistenzsystem} bezeichnet (engl.: \emph{AI Coding Agent}). Der Begriff wird hier im Sinne eines Assistenzsystems ohne eigene Aufgabenplanung verwendet.
In vielen Arbeitsabläufen verschiebt sich der Schwerpunkt von der reinen Codeerzeugung hin zur Auswahl und Integration vorgeschlagener Änderungen.
Damit gewinnt eine Arbeitsweise an Relevanz, die häufig als \emph{Vibe-Coding} beschrieben wird.
Entwicklerinnen und Entwickler delegieren Teilaufgaben an ein KI-Assistenzsystem und übernehmen verstärkt Aufgaben der semantischen Kontrolle, des Reviews sowie der Einbettung in bestehende Strukturen.
Diese Form der Mixed-Initiative-Zusammenarbeit betont weniger syntaktische Details als vielmehr das Verständnis von Anforderungen, Nebenbedingungen und Systemkontext.

Im Brownfield-Kontext ist diese Entwicklung besonders bedeutsam.
Ein Großteil praktischer Softwareentwicklung findet in bestehenden Codebasen statt, die technische Schulden, heterogene Architekturen, historische Entscheidungen und operative Randbedingungen aufweisen.
Bei der Integration von KI-Assistenzsystemen in etablierte Entwicklungsprozesse werden neben Effizienzgewinnen auch Qualitätssicherung, Nachvollziehbarkeit und Risikobegrenzung als relevante Aspekte diskutiert.
Daraus ergibt sich ein Spannungsfeld zwischen schneller Iteration und belastbarem Abgleich.

Für die vorliegende Arbeit ist dabei zentral: Potenzielle Effizienzgewinne werden nicht quantifiziert; im Fokus steht vielmehr, wie Integration, Abgleich und Review im Zusammenspiel mit einem KI-Assistenzsystem empirisch beobachtbar werden.

Wissenschaftlich leistet die Arbeit damit einen Beitrag zur empirischen Auswertung von KI-Assistenzsystemen im Brownfield-Kontext, der über reine Produktivitätsbetrachtungen hinausgeht.

Vor diesem Hintergrund ist der Nutzen eines KI-Assistenzsystems in dieser Arbeit nicht als reine Eigenschaft des Tools zu verstehen, sondern als Zusammenspiel aus (i) Interaktion zwischen Mensch und System, (ii) entwicklerseitiger Kompetenz und (iii) Praktiken des Abgleichs und der Integration.
Automation Bias kann dazu beitragen, dass Vorschläge überbewertet oder unzureichend kontrolliert werden, während umgekehrt übermäßige Skepsis potenzielle Nutzenaspekte reduziert.
Eine wissenschaftlich belastbare Einordnung macht daher eine Betrachtung erforderlich.
Diese verknüpft Kompetenz als Analyseperspektive mit beobachtbaren Artefakten und Ergebnissignalen der Arbeit am Code, ohne daraus kausale Effekte oder Tool-Rangordnungen abzuleiten.

\section{Problemstellung}
In der Diskussion um KI-Assistenzsysteme werden häufig Produktivitätsannahmen formuliert, ohne hinreichend zu differenzieren, wie stark die Interaktion durch Kompetenz, Arbeitsstil und Abgleichspraktiken moderiert wird.
Unter realitätsnahen Bedingungen bestehender Codebasen und unter Zeitdruck bleibt zudem unklar, in welchem Verhältnis wahrgenommener Fortschritt, implementierte Änderungen und als Execution observed ausgewiesenes Verhalten zueinander stehen.
Damit entsteht eine Forschungslücke: Kontrollierte Untersuchungen kombinieren selten Brownfield-Settings, eine differenzierte Ergebnislogik entlang einer Pipeline und Kompetenz als explizite Analyseperspektive.

Diese Arbeit ist daher keine Produktivitätsmessung (z.B. Output-Geschwindigkeit, Codeumfang oder Zeitersparnis), sondern eine Analyse von Integrations-, Abgleichs- und Reviewarbeit im Vibe-Coding-Kontext.

Sie trifft keine Aussagen über kausale Effekte und leitet keine Rangordnung von Tools oder Modellen ab.

Eine zentrale methodische Herausforderung liegt in der Messbarkeit von Zielerreichung.
Chat-Interaktionen und Planäußerungen im Assistenzdialog sind nicht gleichzusetzen mit funktionsfähiger Implementierung.
Ebenso ist eine Implementierung im Code nicht gleichzusetzen mit einem Execution observed-Ausweis dafür, dass gewünschtes Verhalten tatsächlich erreicht wurde.
Diese Differenzen sind im Kontext von KI-generiertem Code besonders relevant, weil die Geschwindigkeit der Generierung und die Plausibilität der Vorschläge die Notwendigkeit eines getrennten Statusausweises nicht aufheben, sondern eher erhöhen können.

Zur Bearbeitung dieser Messprobleme wird eine Diagnoselogik benötigt, die Ergebnisse schrittweise differenziert und durch mehrere Evidenzquellen absichert.
In dieser Arbeit wird dafür eine Pipeline-Logik verwendet, die den Bearbeitungsstand einer Aufgabe als \textit{Attempted}, \textit{Implemented\textsubscript{static}} oder \textit{Execution observed} beschreibt.
Für die Trichterdarstellung wird in Kapitel~6 zusätzlich eine abgeleitete Evidenzstufe  \textit{Implemented}\,$^\ast$ verwendet.
Diese Status sind als Diagnose- und Ergebnislogik zu verstehen, nicht als Qualitätsurteil.
Die Einordnung wird nicht aus einer einzelnen Quelle abgeleitet, sondern trianguliert über Code-Artefakte, Ausführungsartefakte, Chat-Logs und Survey-Daten.
Ziel ist eine methodisch robuste, deskriptive Einordnung, die subjektive oder dialogbasierte Signale nicht mit Execution observed-Ausweisen gleichsetzt und den Interaktionskontext dennoch berücksichtigt.

\section{Zielsetzung der Arbeit}
Ziel dieser Arbeit ist die Auswertung eines KI-Assistenzsystems im Vibe-Coding-Kontext in einer kontrollierten Brownfield-Sandbox.
Zur begrifflichen und methodischen Trennung wird zwischen der Assistenzumgebung GitHub Copilot (Tool), dem zugrundeliegenden LLM Claude Sonnet~4.5 und dem KI-Assistenzsystem als funktionaler Einheit unterschieden.
Beobachtungen beziehen sich je nach Datenquelle auf unterschiedliche Ebenen.
Ein Tool- oder Modellvergleich ist nicht Gegenstand der Arbeit.
Die Wahl von GitHub Copilot erfolgt aus Gründen der ökologischen Validität: Der Integrationskontext (BuyIn) sowie die Zielsetzung der Arbeit orientieren sich an einer Microsoft-nahen Entwicklungsumgebung (insbesondere VS Code und GitHub-Tooling), in der Copilot als praxisnahe Assistenzumgebung naheliegt.
Die Datenerhebung erfolgt in einer kontrollierten Coding-Challenge mit $N = 18$ Teilnehmenden, einer Timebox von 60 Minuten und einer standardisierten, containerisierten Arbeitsumgebung.
Die Aufgaben umfassen acht Tasks (T1--T8).
In der Ergebnisanalyse werden insbesondere frühe Tasks (T1--T5) berücksichtigt, da das Zeitlimit die Bearbeitungsreihenfolge und -tiefe strukturell beeinflusst.

Gegenstand der Arbeit ist dabei ein \emph{interaktives} KI-Assistenzsystem mit kontinuierlicher menschlicher Kontrolle (Human-in-the-Loop); nicht betrachtet werden autonome Agenten mit selbstständiger Aufgabenplanung, Tool-Auswahl oder Entscheidungsautonomie.

Die Arbeit leistet im Kern vier Beiträge:
\begin{enumerate}
    \item[(a)] Entwicklung und Verwendung einer Brownfield-Sandbox als Messumgebung für kontrollierte, aber realitätsnahe Aufgaben in einer bestehenden Codebasis.
    \item[(b)] Operationalisierung von Aufgabenerfolg über die Pipeline-Status: \\
          	extit{Attempted}, \textit{Implemented\textsubscript{static}} und \textit{Execution observed} (sowie \textit{Implemented}\,$^\ast$ als abgeleitete Evidenzstufe) als Trichter- und Diagnoselogik.
    \item[(c)] Analyse der Ergebnisse im Kontext von Kompetenzgruppen als Analyseperspektiven (No-Code, Low-Code, Mid-Code, High-Coder), ohne diese Einteilung als Wertung von Personen zu interpretieren.
    \item[(d)] Methodische Triangulation über Code-Artefakte, Ausführungsartefakte, Chat-Logs und Survey, um die Differenz zwischen intendiertem Vorgehen, implementierten Änderungen und als Execution observed ausgewiesenem Verhalten empirisch einordnen zu können.
\end{enumerate}

\paragraph{Forschungsfragen.}
Aus diesen Zielsetzungen ergeben sich drei Forschungsfragen:
\begin{itemize}
    \item  \textbf{RQ1:} Wie hängt die Entwicklerkompetenz als Analyseperspektive mit dem Erreichen der Pipeline-Status  \textit{Attempted},  \textit{Implemented\textsubscript{static}} und  \textit{Execution observed} (sowie der abgeleiteten Evidenzstufe  \textit{Implemented$^\ast$}) in einer kontrollierten Brownfield-Sandbox zusammen?
    \item  \textbf{RQ2:} Wie tragen Code- und Ausführungsartefakte zur Ergebniseinordnung entlang der Pipeline-Status bei, insbesondere zur Trennung von Implementierung und Execution observed-Ausweisen?
    \item  \textbf{RQ3:} Welche Chat-Interaktionsmuster treten im Vibe-Coding-Kontext auf, und wie sind sie für die Auswertung eines KI-Assistenzsystems im Hinblick auf Abgleich und Integration einzuordnen?
\end{itemize}
Die Kompetenzgruppen dienen dabei ausschließlich der analytischen Typisierung; daraus wird keine Rangordnung abgeleitet und die Einleitung nimmt keine Ergebnisse vorweg.
Die Aussagen dieser Arbeit beziehen sich auf Branches als Analyseeinheiten und auf aggregierte Muster über Gruppen hinweg, nicht auf eine Leistungsbewertung einzelner Personen.
Ergänzend werden theoriegeleitete Perspektiven formuliert, um erwartbare Muster im Spannungsfeld zwischen Unterstützungspotenzial, Abgleichsaufwand und Interaktionsdynamik strukturiert zu beschreiben, ohne Ergebnisse vorwegzunehmen.

\section{Aufbau der Arbeit}
Kapitel~2 führt die theoretischen Grundlagen ein und ordnet Vibe-Coding als Arbeitsstil in relevante Perspektiven ein, darunter Kompetenzmodelle, Human-Computer Interaction, Cognitive Load, Automation Bias sowie etablierte Bezugsrahmen aus Qualitäts- und Produktivitätsdiskursen (u.\,a. ISO/SPACE).
Kapitel~3 beschreibt das Forschungsdesign, die Herleitung theoriegeleiteter Perspektiven und die Operationalisierung zentraler Konstrukte, einschließlich der Kompetenzgruppen als Analyseperspektiven und der Pipeline-Status  \textit{Attempted},  \textit{Implemented\textsubscript{static}} und  \textit{Execution observed} (sowie  \textit{Implemented$^\ast$}).
Kapitel~4 erläutert die Durchführung der Studie, die Datenerhebung und die verwendeten Metriken sowie die Diskussion der Gütekriterien durch die Kombination mehrerer Datenquellen.

Kapitel~5 dokumentiert die technische Brownfield-Sandbox, ihre Architektur und die Automatisierungs- und Auswertungsinfrastruktur, die eine standardisierte, reproduzierbare Messung im kontrollierten Setting ermöglicht.
Kapitel~6 präsentiert die Ergebnisse in deskriptiver Form und strukturiert sie entlang der Pipeline-Logik und der triangulierten Evidenz aus Artefakten und Interaktionen.
Kapitel~7 diskutiert die Befunde, ordnet sie in den theoretischen Kontext ein und fasst Limitationen zusammen, ohne kausale Schlussfolgerungen über die beobachteten Zusammenhänge hinaus zu beanspruchen.
Kapitel~8 schließt die Arbeit mit einer synthetischen Beantwortung der Forschungsfragen und einer Einordnung der Beiträge ab.

\paragraph{Abgrenzung und wissenschaftlicher Beitrag.}
Im Unterschied zu Teilen der bestehenden Forschung steht in dieser Arbeit weder die Optimierung von Prompts noch der Vergleich von Tools oder Modellen im Vordergrund.
Ebenso werden keine Produktivitätsmetriken im Sinne von Ausstoß, Geschwindigkeit oder Zeitersparnis als primäre Zielgrößen herangezogen.
Stattdessen wird Vibe-Coding im Brownfield-Setting als Interaktions- und Integrationsprozess untersucht, der sich nicht zuverlässig über einzelne Proxy-Metriken abbilden lässt.
Der Beitrag liegt in einer auswertenden Perspektive auf Nutzungstypen (Kompetenzgruppen als Analyseperspektiven), einer evidenzbasierten Ergebnislogik über Pipeline-Status sowie der systematischen Triangulation über Code-, Ausführungs-, Chat- und Survey-Artefakte.
Damit wird nachvollziehbar, wie sich intendiertes Vorgehen, implementierte Änderungen und als Execution observed ausgewiesenes Verhalten unter Zeitdruck zueinander verhalten, ohne daraus eine Leistungsbewertung einzelner Personen oder eine Rangordnung von Assistenzsystemen abzuleiten.

\section{Hintergrund und Stand der Forschung}
\label{sec:background}

In diesem Kapitel werden die zentralen Konzepte und theoretischen Grundlagen eingeführt, die für die Bewertung von \emph{AI Coding Agents} im Rahmen von Vibe-Coding in der Unternehmenssoftwareentwicklung relevant sind. Zunächst werden AI Coding Agents definiert und von generischen LLM-basierten Chatbots abgegrenzt (Abschnitt~\ref{sec:definition-agents}). Danach werden Modelle zu Softwarequalität und Produktivität erläutert (Abschnitt~\ref{sec:quality-productivity}), gefolgt von zentralen Konzepten der Human-AI-Collaboration und der Usability (Abschnitt~\ref{sec:human-ai}). Abschnitt~\ref{sec:governance} behandelt Enterprise-Governance-Aspekte wie Sicherheit, Datenschutz und Auditierbarkeit. Abschließend gibt Abschnitt~\ref{sec:related-work} einen Überblick über den Stand der Forschung und verwandte Arbeiten zu konkreten Werkzeugen wie GitHub Copilot, Cursor, Codeium, Plandex und Google Cloud Code/Duet~AI.

\subsection{Definition: AI Coding Agents}
\label{sec:definition-agents}

\emph{AI Coding Agents} sind spezialisierte KI-gestützte Werkzeuge, die Softwareentwicklerinnen und -entwickler bei typischen Programmieraufgaben unterstützen, indem sie auf Basis von Quelltext, Projektkontext und natürsprachlichen Eingaben (\emph{Prompts}) Codevorschläge generieren, bestehenden Code analysieren und teilweise eigenständig Änderungen durchführen. Technisch basieren sie überwiegend auf \emph{Large Language Models (LLMs)}, die auf umfangreichen Mengen von Quelltext und technischer Dokumentation trainiert wurden.

Wesentlich ist die Abgrenzung zu generischen LLM-basierten \emph{Chatbots} (z.\,B. allgemein nutzbare Chat-Assistenten): Ein Chatbot kann zwar ebenfalls Code erzeugen, agiert jedoch primär in einer dialogischen Oberfläche und ist meist nicht tief in Entwicklungswerkzeuge integriert. AI Coding Agents hingegen sind so gestaltet, dass sie direkt in den Entwicklungsworkflow eingebettet sind, etwa als Erweiterung einer integrierten Entwicklungsumgebung (\emph{Integrated Development Environment, IDE}) oder als Kommandozeilen-Werkzeug, das auf das Projektverzeichnis zugreift. Sie sind zudem explizit auf Programmieraufgaben optimiert, etwa durch spezielle Trainingsdaten, Anpassungen der Eingabe-/Ausgabestruktur und Zusatzfunktionen wie Testgenerierung oder Refaktorierung.

In der Praxis lassen sich drei grobe Varianten unterscheiden:

\begin{itemize}
  \item \textbf{IDE-gebundene Agenten:} Diese Tools integrieren sich als Plugin oder Erweiterung in bestehende IDEs wie Visual Studio Code, JetBrains-IDE-Familie oder Visual Studio. Beispiele sind GitHub Copilot, Codeium oder Tabnine. Die Interaktion erfolgt typischerweise inline im Editor: Während der Eingabe werden Codezeilen oder ganze Blöcke vorgeschlagen; ergänzend stehen häufig Chat-Panels zur Verfügung, in denen Entwicklerinnen und Entwickler natürlichsprachliche Fragen stellen können. 
  \item \textbf{CLI-basierte Agenten:} Hier erfolgt die Interaktion über eine \emph{Command Line Interface} (Kommandozeile). Ein Beispiel ist Plandex, ein terminalbasierter Agent, der umfangreiche Aufgaben in Form von mehrstufigen Plänen abarbeitet. Entwickler formulieren Aufgaben auf hoher Abstraktionsebene, woraufhin der Agent Änderungen in mehreren Dateien erzeugt und diese oftmals als Diffs zur Prüfung präsentiert.
  \item \textbf{Hybride Systeme:} Hybride Systeme kombinieren Chat-Funktionalität, kontextuelle Codebearbeitung und teils eigene Editoroberflächen. Cursor etwa stellt einen AI-first-Codeeditor dar, in dem ein integrierter Agent auf Projektebene agiert und über einen Chat-ähnlichen Dialog umfangreiche Refaktorierungen oder Featurerealisierungen ausführen kann.
\end{itemize}

Im Kontext dieser Arbeit werden AI Coding Agents als \emph{„KI-Paarprogrammierer“} verstanden: Sie unterstützen nicht nur bei der Vervollständigung einzelner Codezeilen, sondern übernehmen auch Teile der Problemanalyse, der Lösungsstrukturierung und der Dokumentation. Das Konzept des \emph{Vibe-Coding} beschreibt dabei einen Entwicklungsstil, in dem Entwicklerinnen und Entwickler auf höherer Abstraktionsebene in natürlicher Sprache arbeiten, während die KI die detaillierte Umsetzung in ausführbaren Code übernimmt. Die Verantwortung für die fachliche und technische Korrektheit verbleibt jedoch beim Menschen, der die Vorschläge prüft, ausführt und gegebenenfalls verwirft.

\subsection{Softwarequalität \& Produktivität}
\label{sec:quality-productivity}

Die Einführung von AI Coding Agents wirft unmittelbar die Frage auf, wie sich diese Werkzeuge auf die Qualität der erstellten Software sowie auf die Produktivität der beteiligten Personen auswirken. In diesem Abschnitt werden zunächst etablierte Qualitätsmodelle vorgestellt, bevor Konzepte der Entwicklerproduktivität erläutert und empirische Befunde zum Einfluss von AI Coding Agents diskutiert werden.

\subsubsection*{Softwarequalität: ISO/IEC~25010 und Qualitätsdimensionen}

Ein verbreiteter internationaler Standard zur Beschreibung von Softwarequalität ist ISO/IEC~25010:2011. Der Standard definiert acht Hauptmerkmale der Produktqualität: \emph{funktionale Eignung}, \emph{Zuverlässigkeit}, \emph{Performance-Effizienz}, \emph{Kompatibilität}, \emph{Benutzbarkeit}, \emph{Sicherheit}, \emph{Wartbarkeit} und \emph{Übertragbarkeit}.\cite{iso25010} Jedes dieser Merkmale ist weiter in Submerkmale untergliedert (z.\,B. umfasst Benutzbarkeit Aspekte wie Erlernbarkeit, Bedienbarkeit und Zufriedenstellung).

Für den Einsatz von AI Coding Agents sind insbesondere folgende Dimensionen relevant:

\begin{itemize}
  \item \textbf{Wartbarkeit:} KI-generierter Code kann sowohl positiv als auch negativ auf die Wartbarkeit wirken. Positiv, wenn der Agent konsistente Patterns, klare Strukturierung und ausführliche Kommentare erzeugt; negativ, wenn Vorschläge redundant, unübersichtlich oder inkonsistent sind. Studien zeigen, dass Entwickler KI-Vorschläge teilweise als „Boilerplate-Automat“ einsetzen, was zu konsistenteren Standardstrukturen führen kann, zugleich aber auch das Risiko birgt, unverständliche Codefragmente zu erzeugen, wenn der Kontext nicht korrekt erfasst wird.\cite{nadi2022humaicollab}
  \item \textbf{Sicherheit:} Untersuchungen zu GitHub Copilot weisen darauf hin, dass generierter Code sicherheitsrelevante Schwachstellen enthalten kann, wenn der Agent etwa unsichere Bibliotheken oder Patterns vorschlägt.\cite{pearce2022asleep} Unternehmen müssen daher zusätzliche Sicherheitsmechanismen (Code-Reviews, Security-Scans) einplanen.
  \item \textbf{Benutzbarkeit:} Die Benutzbarkeit der durch KI erzeugten Artefakte (z.\,B. Dokumentation, Tests) sowie die Benutzbarkeit des Tools selbst beeinflussen, ob Entwickler:innen die Vorschläge akzeptieren und effizient nutzen können. ISO/IEC~25010 versteht Usability als Grad, zu dem ein Produkt von bestimmten Benutzern in einem bestimmten Kontext effektiv, effizient und zufriedenstellend verwendet werden kann.\cite{iso25010}
\end{itemize}

Damit wird deutlich, dass eine Bewertung von AI Coding Agents nicht nur auf die syntaktische Korrektheit der generierten Lösungen abstellen darf, sondern auch auf deren Auswirkungen auf Wartbarkeit, Sicherheit und Usability.

\subsubsection*{Produktivität: Begriffe und Messansätze}

Produktivität in der Softwareentwicklung lässt sich allgemein als Verhältnis von Output (z.\,B. implementierte Features, gelöste Tickets) zu Input (v.\,a. eingesetzte Zeit und Personalressourcen) auffassen. Klassische Kennzahlen wie „Lines of Code pro Tag“ wurden in der Literatur umfangreich kritisiert, da sie Qualitätsaspekte und Kontextfaktoren ausblenden.\cite{forsgren2021space} 

Aktuelle Ansätze betonen daher einen multidimensionalen Produktivitätsbegriff. Das \emph{SPACE-Framework} von Forsgren et al.\ beschreibt fünf Dimensionen: \emph{Satisfaction and Well-being}, \emph{Performance}, \emph{Activity}, \emph{Communication and Collaboration} sowie \emph{Efficiency and Flow}.\cite{forsgren2021space} Für AI Coding Agents ist insbesondere relevant, inwieweit sie:

\begin{itemize}
  \item die Bearbeitungszeit für typische Aufgaben reduzieren (\emph{Efficiency}),
  \item die wahrgenommene Arbeitsbelastung senken und damit Wohlbefinden und Zufriedenheit steigern (\emph{Satisfaction}),
  \item die Zusammenarbeit im Team beeinflussen, etwa indem mehr Zeit für konzeptionelle Diskussionen bleibt (\emph{Communication and Collaboration}).
\end{itemize}

Mehrere empirische Studien deuten darauf hin, dass AI Coding Agents die Aufgabendauer signifikant reduzieren können. So zeigte eine kontrollierte Studie, dass Entwickler:innen mit GitHub Copilot bestimmte Programmieraufgaben deutlich schneller bearbeiten konnten als eine Vergleichsgruppe ohne KI-Unterstützung.\cite{peng2023copilot} Eine groß angelegte Feldstudie in Industrieunternehmen berichtete, dass Entwickler:innen mit kontinuierlichem Zugang zu Copilot im Mittel mehr Aufgaben pro Zeiteinheit abschließen konnten, ohne dass sich die Fehlerdichte in der Produktion erhöhte.\cite{cui2025fieldstudy} Gleichzeitig betonen die Autor:innen, dass solcher Produktivitätsgewinn nur dann nachhaltig ist, wenn bestehende Qualitätspraktiken (Tests, Reviews) beibehalten werden.

Subjektive Befragungen unter Entwickler:innen ergänzen diese Befunde: Viele berichten, dass KI-Tools vor allem Routineaufgaben (Boilerplate, einfache Datenzugriffe, Standard-Tests) beschleunigen und monotone Tätigkeiten verringern.\cite{bakal2024zoominfo} Dadurch bleibt mehr Zeit für konzeptionelle Arbeit und Architekturentscheidungen, was wiederum positive Effekte auf Motivation und wahrgenommene Autonomie haben kann.

\subsection{Human--AI Collaboration \& Usability}
\label{sec:human-ai}

Die Wirksamkeit von AI Coding Agents hängt nicht allein von der Modellqualität ab, sondern maßgeblich von der Art und Weise, wie Menschen mit diesen Werkzeugen interagieren. In der Forschung zu Human--AI Collaboration, Usability und Mensch–Automations-Interaktion sind verschiedene theoretische Modelle etabliert, die auch für die Bewertung von Coding Agents relevant sind.

\subsubsection*{Technology Acceptance Model (TAM)}

Das \emph{Technology Acceptance Model (TAM)} nach Davis erklärt die Nutzungsbereitschaft neuer Technologien im Wesentlichen über zwei Wahrnehmungsdimensionen: die \emph{wahrgenommene Nützlichkeit} (\emph{Perceived Usefulness}) und die \emph{wahrgenommene Einfachheit der Nutzung} (\emph{Perceived Ease of Use}).\cite{davis1989tam} Übertragen auf AI Coding Agents bedeutet dies: 

\begin{itemize}
  \item Je stärker Entwickler:innen das Gefühl haben, dass ein KI-Assistent ihre Arbeit spürbar verbessert (z.\,B. durch Zeitersparnis, Fehlerreduktion, bessere Codequalität), desto höher ist die Nutzungsbereitschaft.
  \item Je einfacher das Tool in bestehende Arbeitsabläufe integrierbar ist (z.\,B. reibungslose IDE-Integration, intuitive Shortcuts, geringe Einarbeitungszeit), desto geringer sind Nutzungsbarrieren.
\end{itemize}

Erweiterte TAM-Varianten (z.\,B. TAM2 oder UTAUT) berücksichtigen zusätzlich soziale Einflüsse und organisatorische Rahmenbedingungen. In Teams, in denen der Einsatz von AI Coding Agents aktiv empfohlen wird oder positive Erfahrungsberichte kursieren, steigt typischerweise die Bereitschaft, diese Werkzeuge auszuprobieren und langfristig zu nutzen.

\subsubsection*{Cognitive Load Theory und kognitive Belastung}

Die \emph{Cognitive Load Theory} (CLT) unterscheidet zwischen intrinsischer, extrinsischer und lernbezogener (\emph{germane}) kognitiver Last.\cite{sweller1988clt} AI Coding Agents sollten idealerweise die \emph{extrinsische} Last verringern, etwa indem sie lästige Detailarbeit (Boilerplate, repetitive Strukturen, einfache Transformationsaufgaben) übernehmen und so die mentalen Ressourcen der Entwickler:innen für die eigentliche Problemlösung (intrinsische Last) freisetzen.

Empirische Studien und Erfahrungsberichte deuten darauf hin, dass Entwickler:innen die Arbeit mit KI-Assistenz häufig als entlastend erleben, sofern die Vorschläge relevant und zuverlässig sind.\cite{peng2023copilot,bakal2024zoominfo} Gleichzeitig kann eine schlecht gestaltete Interaktion die kognitive Last erhöhen: Wenn Vorschläge häufig unpassend sind, die Bedienung kompliziert ist oder die KI unerwartete Veränderungen vornimmt, entsteht zusätzlicher Prüf- und Korrekturaufwand. Für die Gestaltung von Vibe-Coding-Szenarien ist daher wichtig, Nutzeroberflächen und Interaktionsmuster so zu entwerfen, dass sie kognitive Belastung minimieren (z.\,B. durch klare Vorschlagsvisualisierung, einfache Möglichkeiten zum Akzeptieren/Ablehnen und konsistente Reaktionen des Systems).

\subsubsection*{Vertrauen in Automation}

In der Forschung zur Mensch–Automations-Interaktion wird \emph{Vertrauen} als zentrale Größe betrachtet.\cite{leesee2004trust,hoffman2018trust} Ein angemessenes Vertrauen in ein automatisiertes System bedeutet, dass Nutzer weder blind vertrauen (\emph{Overtrust}) noch die Automation grundsätzlich ignorieren (\emph{Distrust}), sondern die Fähigkeiten und Grenzen realistisch einschätzen.

Auf AI Coding Agents übertragen heißt dies:

\begin{itemize}
  \item \textbf{Zu hohes Vertrauen} kann dazu führen, dass Entwickler:innen KI-Vorschläge unkritisch übernehmen und etwa Sicherheitsaspekte oder Randfälle nicht ausreichend prüfen.
  \item \textbf{Zu geringes Vertrauen} führt dazu, dass das Tool kaum genutzt wird und sein Potenzial ungenutzt bleibt.
\end{itemize}

Vertrauen wird durch Faktoren wie Transparenz, Zuverlässigkeit und Vorhersagbarkeit geprägt. In Studien zu Copilot und ähnlichen Tools zeigt sich, dass Entwickler:innen das Vertrauen in das Tool schrittweise anpassen: Positive Erfahrungen (z.\,B. gelungene Implementierungen, verlässliche Vorschläge) erhöhen das Vertrauen, während auffällige Fehler oder „zu selbstbewusste“ falsche Antworten es senken.\cite{nadi2022humaicollab} Für die Gestaltung von AI Coding Agents bedeutet dies, dass Mechanismen zur Vertrauenskalibrierung sinnvoll sind, etwa durch Anzeige von Unsicherheiten, klare Grenzen des Einsatzbereichs und die einfache Möglichkeit, Vorschläge zurückzuweisen.

\subsubsection*{Usability eines AI Coding Agents}

Neben den kognitiven und sozialen Aspekten spielt die allgemeine Usability des Tools eine zentrale Rolle. Aufbauend auf ISO/IEC~25010 umfasst Usability u.\,a. Erlernbarkeit, Effizienz, Fehlerrobustheit und Zufriedenstellung.\cite{iso25010} Für AI Coding Agents sind insbesondere folgende Fragen relevant:

\begin{itemize}
  \item Wie leicht ist das Tool zu installieren und zu konfigurieren (z.\,B. IDE-Plugin, Authentifizierung, Projektkonfiguration)?
  \item Wie intuitiv ist die Interaktion (Shortcuts, UI-Design, Sichtbarkeit der Vorschläge)?
  \item In welchem Maß stört oder unterstützt das Tool den \emph{Flow} der Entwickler:innen?
\end{itemize}

Feldberichte, etwa aus der Einführung von Copilot in Unternehmen, zeigen, dass Entwickler:innen das Tool insbesondere dann positiv bewerten, wenn die Integration in die bestehende Tool-Landschaft reibungslos verläuft und die Vorschläge kontextuell passend sind.\cite{bakal2024zoominfo} Übermäßige oder irrelevante Autovervollständigungen werden hingegen als störend empfunden und können dazu führen, dass die Funktion deaktiviert wird.

\subsection{Enterprise-Governance: Sicherheit, Datenschutz und Integration}
\label{sec:governance}

Während einzelne Entwickler:innen AI Coding Agents relativ frei erproben können, ist der Einsatz in Unternehmensumgebungen durch Governance-Anforderungen geprägt. Zu den wichtigsten Themen gehören Informationssicherheit, Datenschutz (insbesondere DSGVO), Auditierbarkeit und Integration in bestehende Prozesse und Werkzeuge.

\subsubsection*{Sicherheit und Codequalität}

Aus Security-Perspektive stellen AI Coding Agents eine zweifache Herausforderung dar: Einerseits kann der generierte Code selbst Sicherheitslücken enthalten, andererseits kann das Werkzeug zum ungewollten Abfluss sensibler Informationen führen.

Studien zu Copilot zeigen, dass generierte Codefragmente in sicherheitskritischen Szenarien eine erhöhte Wahrscheinlichkeit aufweisen, bekannte Schwachstellen zu enthalten, etwa unsichere Standardkonfigurationen oder fehlerhafte Eingabevalidierung.\cite{pearce2022asleep} Deshalb empfehlen zahlreiche Autoren, KI-generierten Code grundsätzlich denselben Sicherheitsprüfungen (Code-Review, automatisierte Security-Scans, Penetrationstests) zu unterziehen wie manuell geschriebenen Code.\cite{nadi2022humaicollab}

Zudem müssen Unternehmen verhindern, dass vertrauliche Informationen (z.\,B. Geschäftslogik, API-Schlüssel, personenbezogene Daten) über KI-Services nach außen gelangen. Einige Anbieter bieten dafür spezifische Enterprise-Varianten an, in denen zugesichert wird, dass Prompts und Code nicht für Trainingszwecke verwendet und nur kurzfristig verarbeitet werden.\cite{githubcopilottrust} Andere Tools wie Codeium werben damit, dass sie mit selbst-hostbaren Varianten vollständig im Rechenzentrum des Kunden betrieben werden können und Trainingsdaten ausschließlich aus permissiven Open-Source-Projekten stammen, um Lizenz- und Datenschutzrisiken zu minimieren.\cite{ramel2023codeium}

\subsubsection*{Datenschutz und DSGVO-Compliance}

Die europäische Datenschutz-Grundverordnung (DSGVO) verlangt bei der Verarbeitung personenbezogener Daten strenge Regeln hinsichtlich Rechtsgrundlage, Zweckbindung, Transparenz und Datensicherheit. Enthält der übermittel­te Kontext personenbezogene Daten oder lassen sich aus dem Code Rückschlüsse auf Individuen ziehen, fungiert ein Cloud-basierter AI Coding Agent als Auftragsverarbeiter im Sinne der DSGVO. Unternehmen müssen dann prüfen:

\begin{itemize}
  \item In welche Länder werden die Daten übertragen?
  \item Welche Speicherfristen gelten?
  \item Werden Daten zu Trainingszwecken wiederverwendet?
\end{itemize}

Viele Anbieter stellen hierfür Verträge zur Auftragsverarbeitung bereit und heben Zertifizierungen (z.\,B. ISO~27001) hervor. Dennoch bevorzugen einige Organisationen On-Premises- oder Self-Hosted-Lösungen, um maximale Kontrolle über die Datenflüsse zu behalten.\cite{ramel2023codeium}

Ein weiterer Governance-Aspekt betrifft Lizenz-Compliance: Falls KI-Modelle auf nicht-permissiven Open-Source-Lizenzen (z.\,B. GPL) trainiert wurden, kann nicht ausgeschlossen werden, dass generierte Codefragmente lizenzrechtlich problematisch sind.\cite{vaithilingam2022codecompletion} Einige Anbieter adressieren dies explizit, indem sie solche Lizenzen aus den Trainingsdaten ausschließen oder Ausgaben filtern.

\subsubsection*{Auditierbarkeit und Nachvollziehbarkeit}

In regulierten Branchen besteht häufig die Pflicht, Entwicklungsprozesse nachvollziehbar zu dokumentieren. Wenn AI Coding Agents zum Einsatz kommen, stellt sich die Frage, wie deren Beitrag zum Code historisiert werden kann. 

Mögliche Maßnahmen umfassen:

\begin{itemize}
  \item Logging von Prompts und Antworten in internen Systemen,
  \item Kennzeichnung von KI-generiertem Code in Commit-Messages oder Pull-Requests,
  \item spezielle Review-Regeln für Code, der überwiegend aus KI-Vorschlägen stammt.
\end{itemize}

Einige Enterprise-Angebote von AI Coding Agents stellen erweiterte Audit-Logs und Administrationsfunktionen bereit, mit denen Nutzungsstatistiken, Aktivitätsmuster und Konfigurationen zentral verwaltet werden können. Dies dient sowohl der Compliance als auch der Steuerung der Toolnutzung (z.\,B. schrittweise Rollouts, Lizenzverwaltung).

\subsubsection*{Integration in bestehende Toolchains und Prozesse}

Damit AI Coding Agents in einem Unternehmen wirksam eingesetzt werden können, müssen sie in bestehende Toolchains (z.\,B. Versionsverwaltung, CI/CD-Pipelines, Issue-Tracker) integriert werden. Typische Anforderungen sind:

\begin{itemize}
  \item Unterstützung gängiger IDEs und Editor-Umgebungen,
  \item Single-Sign-On (SSO) und zentrale Benutzerverwaltung,
  \item Konfigurierbarkeit hinsichtlich Sprachen, Frameworks und unternehmensspezifischer Konventionen,
  \item Option, interne Codebasen als Wissensquelle für die KI zu nutzen (z.\,B. für projektspezifische Patterns).
\end{itemize}

Beispiele hierfür sind GitHub Copilot~Enterprise und Googles Duet~AI, die sich explizit an Unternehmenskunden richten und Funktionen wie SSO, erweiterte Administration, Telemetrie und (teilweise) Anpassung an unternehmensspezifische Repositories bieten.\cite{githubcopilottrust,googleduetai} Für ein Unternehmen wie BuyIn ist entscheidend, dass ein gewählter AI Coding Agent sich in bereits etablierte Werkzeuge integrieren lässt und Governance-Vorgaben (z.\,B. Sicherheitsrichtlinien, Freigabeprozesse) nicht unterläuft, sondern idealerweise unterstützt.

\subsection{Related Work: KI-gestützte Softwareentwicklung und aktuelle Tools}
\label{sec:related-work}

In den letzten Jahren ist eine wachsende Zahl wissenschaftlicher Publikationen und Praxisberichte zur KI-unterstützten Softwareentwicklung erschienen. Dieser Abschnitt gibt einen Überblick über zentrale Arbeiten mit Bezug auf AI Coding Agents und Vibe-Coding.

\subsubsection*{Empirische Studien zu GitHub Copilot}

GitHub Copilot als eines der ersten und bekanntesten Produkte in diesem Bereich ist Gegenstand mehrerer empirischer Untersuchungen. \textcite{peng2023copilot} zeigen in einem kontrollierten Experiment, dass Entwickler:innen mit Copilot deutlich schneller Aufgaben bearbeiten können als ohne, wobei die Qualität der Lösungen weitgehend vergleichbar bleibt. \textcite{cui2025fieldstudy} untersuchen Copilot in einer groß angelegten Feldstudie in mehreren Unternehmen und berichten von signifikanten Produktivitätsgewinnen, insbesondere bei weniger erfahrenen Entwickler:innen. \textcite{pearce2022asleep} analysieren demgegenüber die Sicherheitsaspekte und finden, dass ein nicht unerheblicher Anteil der von Copilot vorgeschlagenen Lösungen in sicherheitskritischen Szenarien Schwachstellen enthält.

\textcite{vaithilingam2022codecompletion} sowie \textcite{nadi2022humaicollab} untersuchen die Interaktion zwischen Entwickler:innen und modernen Codevervollständigungswerkzeugen und beschreiben Muster von Vertrauen, Misstrauen und Strategien zur Fehlerkorrektur. Sie zeigen u.\,a., dass Entwickler:innen KI-Vorschläge eher akzeptieren, wenn sie als Ergänzung zu bestehendem Code dienen, während Vorschläge, die größere Strukturänderungen betreffen, skeptischer betrachtet werden.

\subsubsection*{Praxisberichte und Fallstudien}

Neben den klassischen Forschungspublikationen gewinnt „graue Literatur“ in Form von technischen Whitepapern und Unternehmensberichten an Bedeutung. GitHub veröffentlichte mehrere Berichte zur Einführung von Copilot in Organisationen, die von hoher wahrgenommener Produktivität und Zufriedenheit berichten, aber auch Grenzen und Risiken benennen.\cite{githubcopilottrust} \textcite{bakal2024zoominfo} beschreiben etwa die Einführung von Copilot bei ZoomInfo mit mehreren hundert Entwickler:innen. Sie dokumentieren Annahmequoten von rund einem Drittel der Vorschläge und schätzen, dass dadurch etwa 20\,\% der Entwicklungszeit eingespart werden konnten.

Codeium positioniert sich in Whitepapern als datenschutzfreundliche Alternative zu Copilot, die sich besonders für Unternehmen eignet: Zum einen durch den Verzicht auf nicht-permissive Open-Source-Lizenzen im Training, zum anderen durch die Möglichkeit des Self-Hostings.\cite{ramel2023codeium} Plandex wird in Entwicklerblogs als experimenteller Agent für große Aufgaben beschrieben, der projek­tweite Änderungen über die Kommandozeile orchestriert, allerdings bislang ohne formale Evaluation in industriellen Kontexten.

\subsubsection*{Human--AI Collaboration in der Softwareentwicklung}

Über konkrete Tools hinaus existiert eine Reihe von Arbeiten, die die Zusammenarbeit zwischen Menschen und KI allgemein untersuchen. \textcite{nadi2022humaicollab} geben einen Überblick über Patterns der Mensch--KI-Kollaboration in der Softwareentwicklung und identifizieren typische Herausforderungen wie mangelnde Transparenz, Vertrauen und Verantwortungsteilung. \textcite{forsgren2021space} und andere Produktivitätsstudien liefern theoretische Rahmen, um Effekte von KI-Werkzeugen über reine Zeitmessungen hinaus einzuordnen.

Insgesamt lässt sich festhalten, dass AI Coding Agents erhebliches Potenzial besitzen, die Produktivität und Entwicklererfahrung zu verbessern, gleichzeitig aber offene Fragen hinsichtlich Qualität, Sicherheit, Governance und langfristiger Kompetenzentwicklung aufwerfen. Die vorliegende Arbeit knüpft an diese Forschung an, indem sie eine vergleichende Bewertung mehrerer AI Coding Agents im Kontext von Vibe-Coding und Unternehmenssoftwareentwicklung vornimmt und eine Integrationsstrategie am Beispiel von BuyIn ableitet.

\chapter{Forschungsdesign und theoriegeleitete Perspektiven}
\label{chap:forschungsdesign}

Dieses Kapitel leitet das empirische Design der Studie aus den in Kapitel~\ref{chap:hintergrund} vorgestellten Theorien ab. Ziel ist es, den Zusammenhang der unabhängigen Variable \emph{Kompetenz} mit den beobachteten Dimensionen \emph{Aufgabenfortschritt/Statusausweise} (artefaktgebunden sowie automatisiert erfasst) sowie \emph{Interaktion} in einer kontrollierten Brownfield-Umgebung zu beschreiben.

\section{Methodischer Ansatz: Between-Subjects Design}

Die Untersuchung folgt einem \emph{Between-Subjects Design}, bei dem Teilnehmende basierend\linebreak
auf ihrer Vorkenntnis in disjunkte Gruppen eingeteilt werden. Begründung nach\linebreak
    \citeauthor{wohlin2012experimentation} (\citeyear{wohlin2012experimentation}):
\begin{itemize}
    \item Vermeidung von Lerneffekten (Carry-Over Effects): Ein Within-Subjects\linebreak
    Design (jede Person löst Aufgaben einmal mit, einmal ohne KI) wäre hier proble-\linebreak
    matisch, da das Lösen einer Aufgabe Wissen für den zweiten Durchgang generiert.
    \item Isolierung der Kompetenzvariable: Das Ziel ist der Vergleich der Interaktions-\linebreak
    muster zwischen verschiedenen Kompetenzniveaus (\enquote{Mensch A + Maschine} vs.\linebreak
    \enquote{Mensch B + Maschine}), nicht der Vergleich des Tools an sich.
\end{itemize}

\section{Theoriegeleitete Perspektiven}

Die Perspektiven werden direkt aus dem Dreyfus-Modell, der Dual-Process Theory und der Cognitive Load Theory (siehe Abschnitt~\ref{sec:competence-models} und \ref{sec:human-ai}) abgeleitet.

Zur Orientierung wird die primäre Zuordnung der Perspektiven zu den Forschungsfragen explizit gemacht: P1 adressiert primär RQ1, P2 adressiert primär RQ1 und RQ2, und P3 adressiert primär RQ3. Diese Zuordnung dient ausschlie\ss{}lich der Strukturierung und nimmt keine Ergebnisse vorweg.

\subsection{P1: Der \enquote{Sweet Spot} der KI-Unterstützung (Mid-Code)}
\noindent\csname textbf\endcsname{Theoretische Begründung:} Nach der Cognitive Load Theory profitieren Lernende am meisten, wenn \emph{Extraneous Load} (Syntax) reduziert wird, solange genug Schema-Wissen vorhanden ist, um Output im Kontext einzuordnen. Mid-Code-Teilnehmende befinden sich in einer Phase, in der sie Konzepte verstehen, aber die syntaktische Umsetzung noch kognitive Ressourcen bindet. In dieser Perspektive wird die Assistenzinteraktion als Delegation von Syntaxarbeit verstanden; dadurch kann \emph{Germane Load} für architektonische Entscheidungen freigesetzt werden.
\noindent\csname textbf\endcsname{Alternative Perspektive:} Ein \enquote{Rising Tide}-Narrativ würde nahelegen, dass alle Gruppen gleichermaßen profitieren. Dies passt jedoch nur eingeschränkt zur Theorie der \emph{Zone der nächsten Entwicklung}, da Novizen ohne Scaffolding überfordert sein können.

\begin{quote}
Für die deskriptive Einordnung wird betrachtet, wie häufig hoher Pipeline-Fortschritt innerhalb des 60-Minuten-Limits in den Kompetenzgruppen auftritt (operationalisiert als maximale Anzahl \texttt{implemented}-Markierungen pro Branch in \nolinkurl{evaluation/task_pipeline_v3.json}; \texttt{implemented} wird dabei als heuristisches Artefakt-Signal \textit{Implemented\textsubscript{static}} verstanden).
\end{quote}

\subsection{P2: Das Abgleich-Paradoxon (Low-Code)}
\noindent\csname textbf\endcsname{Theoretische Begründung:} Novizen und Advanced Beginners (Dreyfus Stufe 1--2) fehlt das mentale Modell, um subtile Fehler zu erkennen. Sie operieren primär im \emph{System 1} (intuitiv) und akzeptieren plausibel klingende Vorschläge unkritisch (\emph{Errors of Commission}). Da ihnen die Schemata fehlen, um den Code zu \enquote{parsen}, ist der \emph{Intrinsic Load} des Abgleichs hoch.
\noindent\csname textbf\endcsname{Empirische Bezugnahme:} Studien von \textcite{vaithilingam2022codecompletion} berichten, dass Nutzer oft Code akzeptieren, den sie nicht verstehen, was zu schwer findbaren Bugs führt.

\begin{quote}
Für die deskriptive Einordnung wird betrachtet, wie sich Pipeline-Fortschritt zwischen den Kompetenzgruppen verteilt (operationalisiert über die Anzahl \texttt{implemented}-Markierungen pro Branch in \nolinkurl{evaluation/task_pipeline_v3.json}; \textit{Implemented\textsubscript{static}} als heuristisches Artefakt-Signal). Aussagen zu Sicherheitslücken, Wartbarkeit oder \emph{Overtrust} werden in dieser Studie als \emph{explorativ} behandelt.
\end{quote}

\subsection{P3: Skeptizismus und Qualitätsfokus (High-Coder)}
\noindent\csname textbf\endcsname{Theoretische Begründung:} Experten (Dreyfus Stufe 4--5) haben starke mentale Modelle und vertrauen der Automation weniger. Sie aktivieren \emph{System 2} (analytisch) häufiger, um Vorschläge zu kontrollieren. Dies kann mit einem \emph{Distrust}-Verhalten einhergehen, bei dem potenzielle Zeitersparnisse durch manuelle Reviews teilweise relativiert werden.
\noindent\csname textbf\endcsname{Alternative Sichtweise:} Man könnte annehmen, Experten seien am schnellsten. Die Theorie der \emph{Expertise Reversal Effect} (Kalyuga et al.) legt jedoch nahe, dass Hilfestellungen für Experten redundant und störend sein können.

\begin{quote}
Für die deskriptive Einordnung wird betrachtet, wie häufig maximaler\linebreak
Pipeline\mbox{-}Fortschritt in den Kompetenzgruppen auftritt (operationalisiert über\linebreak
die Anzahl als \texttt{implemented} markierter Tasks pro Branch in \nolinkurl{evaluation/}\linebreak
\nolinkurl{task_pipeline_v3.json}; \textit{Implemented\textsubscript{static}} als heuristisches Artefakt-Signal).
\end{quote}

\begin{quote}
Aussagen zu selektiver Nutzung (z.B. \enquote{Accept All}), Codier-Geschwindigkeit und Codequalität (Wartbarkeit/Sicherheit) werden in dieser Studie als \emph{explorativ} behandelt.
\end{quote}

Die Perspektiven werden in dieser Arbeit theoriegeleitet formuliert und im Sinne einer explorativen Strukturierung der Ergebnisdiskussion verwendet. Eine inferenzstatistische Auswertung erfolgt nicht; die Perspektiven dienen der geordneten Einordnung der in Kapitel~6 berichteten Befundlinien im Kontext der Forschungsfragen. Entsprechend werden die Perspektiven als deskriptive, theoriegeleitete Deutungsfolie unter den dokumentierten Studienbedingungen herangezogen, ohne kausale Effekte oder gruppenstabile Leistungsrangfolgen zu beanspruchen.

\section{Operationalisierung der Kompetenzstufen}

Die Kompetenzgruppen (No-Code/Low-Code/Mid-Code/High-Coder) werden in dieser Arbeit ausschlie\ss{}lich als \emph{Analyseperspektiven} verwendet (Stratifizierung/Segmentierung) und nicht als Leistungsbewertung einzelner Personen. Die Gruppenzuordnung erfolgt auf Basis eines transparenten, auditierbaren Scoring-Modells mit insgesamt 0--13 Punkten. Die Items werden im Pre-Survey erhoben und in Punkte gemappt; die Summe definiert die Kompetenzstufe. Formale Bildungsindikatoren (z.B. Ausbildungsstand/Abschluss) werden bewusst nicht erhoben und nicht verwendet.

\noindent\emph{Hinweis:} Das Scoring stellt eine pragmatische Heuristik dar und ist kein psychometrisch fundiertes Instrument.

\noindent \textbf{Schwellenwerte zur Gruppenzuordnung.} Aus dem Gesamtscore \(S\in[0,13]\) ergeben sich die Kompetenzgruppen wie folgt:  \textbf{No-Code} (\(S=0\)),  \textbf{Low-Code} (\(S=1\) bis \(6\)),  \textbf{Mid-Code} (\(S=7\) bis \(10\)) und  \textbf{High-Coder} (\(S=11\) bis \(13\)). Diese Cutoffs werden \emph{a priori} festgelegt, um nachtr\"agliche Anpassungen an die beobachteten Ergebnisse zu vermeiden.


\noindent \textbf{Reduktion von Self-Report-Bias und Plausibilit\"atschecks.} Da die Operationalisierung auf Selbstauskunft beruht, werden ausschlie\ss{}lich diskrete Erfahrungsangaben (Jahre allgemeine Programmiererfahrung; Jahre TS/Node/Backend) und zwei Likert-Proxies (Implementierungssicherheit; Debugging-/Testpraxis) in ein festes Punkteschema \emph{ohne} nachtr\"agliche Einzelfallkorrekturen \"ubersetzt.

\noindent \textbf{Optionale Sensitivit\"atsanalyse (Future Work).} Als optionale Erweiterung (nicht Teil dieser Arbeit) k\"onnte in zuk\"unftiger Arbeit untersucht werden, wie stabil die beschriebenen Muster gegen plausible Alternativen der Kompetenz-Operationalisierung sind (z.\,B. leicht verschobene Cutoffs oder Variationen der Punktmappings der Selbsteinsch\"atzungsitems). Solche Sensitivit\"atsl\"aufe w\"urden hier lediglich die Abh\"angigkeit der Interpretation von der gew\"ahlten Gruppendefinition transparent machen; es wird dabei keine bestimmte Richtung oder St\"arke eines Effekts erwartet.

\noindent \textbf{Threats to validity (Kompetenzmessung).} Kompetenz ist ein latentes Konstrukt und kann hier nur \emph{approximiert} werden (Konstruktvalidit\"at) \textendash{} \"uber Selbstauskunft zu allgemeiner Programmiererfahrung, TS/Node/Backend-Erfahrung sowie zwei erfahrungsnahen Selbsteinsch\"atzungs-Proxies (Implementierungssicherheit; Debugging-/Testpraxis). Selbstauskunft kann systematisch verzerrt sein (Selbst\"ubersch\"atzung/soziale Erw\"unschtheit), und Likert-Skalen k\"onnen zwischen Personen unterschiedlich interpretiert werden (Messinvarianz). Gegenma\ss{}nahmen sind die a priori festgelegte, auditierbare Rubrik mit diskreten Mappings sowie die als optionale Erweiterung vorgeschlagene Sensitivit\"atsanalyse gegen alternative Cutoffs/Mappings. Die Kompetenzgruppen bleiben damit eine interpretierbare Analyseperspektive, ohne Einzelfallbewertungen abzuleiten.

\section{Kontrolle von Störvariablen (Confounder)}
Um die interne Validität zu sichern, wurden folgende Störvariablen kontrolliert:
\begin{itemize}
    \item  \textbf{IDE-Familiarity:} Alle Teilnehmenden nutzten VS Code. Unterschiede in der Vertrautheit wurden im Pre-Survey erfasst.
    \item  \textbf{Domain Knowledge:} Das Todo-App-Szenario ist generisch genug, um kein spezifisches Domänenwissen (wie z.B. in der Medizin) vorauszusetzen.
    \item  \textbf{Englischkenntnisse:} Da Copilot primär auf englischem Code trainiert ist, wurden Grundkenntnisse vorausgesetzt.
\end{itemize}

\section{Die Sandbox-Umgebung als unabhängige Variable: Ecological Validity}

Die Wahl einer  \textbf{Brownfield-Umgebung} (bestehender Code) statt Greenfield (leeres Blatt) ist essenziell für die \emph{ökologische Validität} der Studie.
Im industriellen Anwendungskontext arbeiten Entwickler:innen zu ca. 90\% in bestehenden Systemen \parencite{brandtner2024brownfield}.
\begin{itemize}
    \item  \textbf{Kontext-Sensitivität (RAG-Szenario):} Im Brownfield-Setting sind bestehende Stil- und Strukturkonventionen (z.\,B. SCSS-Module, React Hooks) relevanter Kontext. Damit werden sowohl Kontexthinweise im Projekt als auch die Praxis adressiert, relevante Dateien zu öffnen und Kontext bereitzustellen.
    \item  \textbf{Kognitive Belastung:} Die Teilnehmenden sind mit der Notwendigkeit konfrontiert, sich in der Struktur zurechtzufinden, was für Low-Code-Teilnehmende eine zusätzliche Hürde darstellt (\emph{Intrinsic Load}). Greenfield-Aufgaben (z.B. LeetCode) würden diese Komplexität ausblenden und die Ergebnisse verzerren.
\end{itemize}
Damit werden \emph{reale} Arbeitsbedingungen simuliert, ohne die Aufgabenbearbeitung auf das Erproben isolierter Algorithmen zu reduzieren.

\begin{table}[ht]
    \centering
    \begingroup
    \footnotesize
    \renewcommand{\arraystretch}{1.25}
    \newcommand{\NoExtract}[1]{\BeginAccSupp{method=escape,ActualText={}}#1\EndAccSupp{}}
    \newcommand{\ExtractOnly}[1]{\smash{\rlap{\BeginAccSupp{unicode,method=pdfstringdef,ActualText={#1}}\textcolor{white}{.}\EndAccSupp{}}}}
    \def\CompetenceRubricActualText{Input-Quelle^^J^^JSkala (Pre-Survey)^^J^^JJahre allgemeine^^JKeine Erfahrung^^JProgrammiererWeniger als 1 Jahr^^Jfahrung^^J1 bis 3 Jahre^^JMehr als 3 Jahre^^JJahre Erfahrung^^Jmit TS/Node/^^JBackend^^J^^JImplementierungs^^Jsicherheit^^J(Selbsteinschätzung)^^JDebugging und^^JTestpraxis^^J(Selbsteinschätzung)^^J^^JKeine Erfahrung^^JWeniger als 1 Jahr^^J1 bis 3 Jahre^^JMehr als 3 Jahre^^JLikert 1–5^^J^^JLikert 1–5^^J^^JPunkte^^J(Mapping)^^J^^J(1)=0^^J(2)=2^^J(3)=4^^J(4)=5^^J(1)=0^^J(2)=1^^J(3)=3^^J(4)=4^^J1–2→0^^J3→1^^J4–5→2^^J^^J1–2→0^^J3→1^^J4–5→2^^J^^JMax. Begründung^^JPunkte^^J5/13^^J^^JProxy für routinierte^^JProgrammierpraxis; als^^JSelbstauskunft erhoben^^Jund daher kein objektives^^JKompetenzmaß.^^J^^J4/13^^J^^JProxy für Domänen- und^^JStack-Nähe (relevant für^^Jdie Brownfield-Sandbox);^^JSelbstauskunft, nicht^^JLeistungsnachweis.^^J^^J2/13^^J^^JProxy für wahrgenommene^^JFähigkeit, Feature-Logik^^Jselbstständig umzusetzen;^^Jsubjektiv und^^Jkontextabhängig.^^J^^J2/13^^J^^JProxy für Validierungsund Fehlersuchroutinen^^J(z.B. Tests/Debugging) im^^JArbeitsprozess;^^JSelbstauskunft statt^^Jobjektiver Messung.^^J^^JTabelle 3.1: Auditierbare Scoring-Rubrik zur Kompetenz-Operationalisierung (0–13 Punkte) basierend auf den tatsächlich erhobenen Pre-Survey-Items (Form Responses 1 in thesis/surveys/Probanden-Einstufungsformular(Responses)^^J.xlsx). Verwendet werden ausschließlich erfahrungsnahe Proxies (Jahre allgemeine Programmiererfahrung; Jahre TS/Node/Backend; Selbsteinschätzung Implementierungssicherheit; Selbsteinschätzung Debugging/Testpraxis). Abk.: TS = TypeScript; Node = Node.js; Backend = Serverteil.^^JDie Gewichtung ist als maximal erreichbare Punktzahl pro Item angegeben;^^Jder Gesamtscore ist die Summe 𝑆.}
    \begin{tabular}{|>{\raggedright\arraybackslash}p{0.17\linewidth}|>{\raggedright\arraybackslash}p{0.28\linewidth}|>{\raggedright\arraybackslash}p{0.14\linewidth}|>{\centering\arraybackslash}p{0.07\linewidth}|>{\raggedright\arraybackslash}p{0.26\linewidth}|}
        \hline
        \NoExtract{{\bfseries Input-Quelle}}\ExtractOnly{\CompetenceRubricActualText} & \NoExtract{{\bfseries Skala (Pre-Survey)}} & \NoExtract{{\bfseries Punkte (Mapping)}} & \NoExtract{{\bfseries Max. Punkte}} & \NoExtract{{\bfseries Begr\"undung}} \\ \hline
        \NoExtract{\parbox[t]{\linewidth}{\vspace{0pt}Jahre allgemeine\\Programmiererfahrung}} & \NoExtract{\parbox[t]{\linewidth}{\vspace{0pt}Keine Erfahrung\par Weniger als 1 Jahr\par 1 bis 3 Jahre\par Mehr als 3 Jahre}} & \NoExtract{\parbox[t]{\linewidth}{\vspace{0pt}(1)=0\par (2)=2\par (3)=4\par (4)=5}} & \NoExtract{5/13} & \NoExtract{Proxy f\"ur routinierte Programmierpraxis; als Selbstauskunft erhoben und daher kein objektives Kompetenzma\ss{}.} \\ \hline
        \NoExtract{Jahre Erfahrung mit TS/Node/Backend} & \NoExtract{\parbox[t]{\linewidth}{\vspace{0pt}Keine Erfahrung\par Weniger als 1 Jahr\par 1 bis 3 Jahre\par Mehr als 3 Jahre}} & \NoExtract{\parbox[t]{\linewidth}{\vspace{0pt}(1)=0\par (2)=1\par (3)=3\par (4)=4}} & \NoExtract{4/13} & \NoExtract{Proxy f\"ur Dom\"anen- und Stack-N\"ahe (relevant f\"ur die Brownfield-Sandbox); Selbstauskunft, nicht Leistungsnachweis.} \\ \hline
        \NoExtract{\parbox[t]{\linewidth}{\vspace{0pt}Implementierungs\-sicherheit\par (Selbst\-einsch\"atzung)}} & \NoExtract{Likert 1--5} & \NoExtract{\parbox[t]{\linewidth}{\vspace{0pt}1--2\(\rightarrow\)0\par 3\(\rightarrow\)1\par 4--5\(\rightarrow\)2}} & \NoExtract{2/13} & \NoExtract{Proxy f\"ur wahrgenommene F\"ahigkeit, Feature-Logik selbstst\"andig umzusetzen; subjektiv und kontextabh\"angig.} \\ \hline
        \NoExtract{\parbox[t]{\linewidth}{\vspace{0pt}Debugging- und Testpraxis\par (Selbst\-einsch\"atzung)}} & \NoExtract{Likert 1--5} & \NoExtract{\parbox[t]{\linewidth}{\vspace{0pt}1--2\(\rightarrow\)0\par 3\(\rightarrow\)1\par 4--5\(\rightarrow\)2}} & \NoExtract{2/13} & \NoExtract{Proxy f\"ur Validierungs- und Fehlersuchroutinen (z.B. Tests/Debugging) im Arbeitsprozess; Selbstauskunft statt objektiver Messung.} \\ \hline
    \end{tabular}

    \NoExtract{\caption[Auditierbare Scoring-Rubrik zur Kompetenz-Operationalisierung]{Auditierbare Scoring-Rubrik zur Kompetenz-Operationalisierung (0--13 Punkte) basierend auf den tats\"achlich erhobenen Pre-Survey-Items (Form Responses 1 in \protect\path{thesis/surveys/Probanden-Einstufungsformular(Responses).xlsx}). Verwendet werden ausschlie\ss{}lich erfahrungsnahe Proxies (Jahre allgemeine Programmiererfahrung; Jahre TS/Node/Backend; Selbsteinsch\"atzung Implementierungssicherheit; Selbsteinsch\"atzung Debugging/Testpraxis). Abk.: TS = TypeScript; Node = Node.js; Backend = Serverteil. Die Gewichtung ist als maximal erreichbare Punktzahl pro Item angegeben; der Gesamtscore ist die Summe \(S\).}}
    \label{tab:competence-rubric}
    \endgroup
\end{table}

\begin{table}[ht]
    \centering
    \begin{tabular}{|l|l|l|>{\raggedright\arraybackslash}p{6cm}|}
    \hline
            \csname textbf\endcsname{Gruppe} &  \textbf{Dreyfus-Level} &  \textbf{Score} &  \textbf{Charakteristik} \\ \hline
        No-Code & Novice & 0 & Kein mentales Modell von Programmierung. Stützt sich stark auf KI als Black Box. \\ \hline
        Low-Code & Advanced Beginner & 1--6 & Kennt Grundkonzepte, aber keine Architekturmuster. Erhöhtes Risiko für \emph{Spaghetti-Code} und \emph{Automation Bias}. \\ \hline
        Mid-Code & Competent & 7--10 & Versteht Zusammenhänge, plant Schritte. Fehler in Vorschlägen sind in dieser Perspektive häufig erkennbar (\emph{System 2} Aktivierung möglich). \\ \hline
        High-Coder & Proficient/Expert & 11--13 & Intuitives Verst\"andnis. Nutzt KI zur Delegation von Boilerplate, nicht zur L\"osungsfindung. \\ \hline
    \end{tabular}
    \caption{Mapping von Kompetenzgruppen auf das Dreyfus-Modell (konzeptionelle, theoriegeleitete Einordnung; idealtypische Kurzcharakteristiken ohne individuelle Leistungszuschreibung).}
    \label{tab:competence-mapping}
\end{table}

\chapter{Durchführung und Methodik der Datenerhebung}
\label{chap:methodik}

Während das vorangegangene Kapitel das Forschungsdesign und die Hypothesen definiert hat, beschreibt dieses Kapitel den konkreten Ablauf der Studie sowie die Operationalisierung der Messgrößen. Es wird detailliert dargelegt, wie die Datenerhebung organisiert wurde, um die interne Validität der Ergebnisse sicherzustellen.

\section{Ablauf der Untersuchung}

Die Studie wurde als kontrolliertes Experiment durchgeführt. Jede Sitzung folgte einem strikten Protokoll, um Störfaktoren zu minimieren. Der Ablauf gliederte sich in fünf Phasen:

\begin{enumerate}
    \item  \textbf{Onboarding (10 Min.):}
    Die Teilnehmenden erhielten eine Einführung in das Szenario (\enquote{Sie sind neuer Entwickler im Team...}) und die Entwicklungsumgebung. Es wurde sichergestellt, dass GitHub Copilot aktiviert und funktionsfähig war.
    
    \item  \textbf{Pre-Survey (5 Min.):}
    Erfassung der demografischen Daten und Selbsteinschätzung zur Kompetenzeinstufung (siehe Kapitel~\ref{chap:forschungsdesign}).
    
    \item  \textbf{Vibe-Coding-Challenge (60 Min.):}
    Die Kernphase der Untersuchung. Die Teilnehmenden bearbeiteten die Aufgaben selbstständig. Eingriffe durch den Versuchsleiter erfolgten nur bei technischen Problemen (z.\,B. Absturz der IDE), nicht jedoch bei inhaltlichen Fragen.
	\csname textbf\endcsname{Begründung des Zeitlimits:} Die Begrenzung auf 60 Minuten dient der organisatorischen Machbarkeit und fungiert als Stressor. Zeitdruck erhöht die Wahrscheinlichkeit, dass Teilnehmende auf heuristische Strategien (\emph{System 1}) zurückgreifen. Dies kann den \emph{Automation Bias} verstärken und simuliert realistische Deadline-Situationen.
    
    \item  \textbf{Datensicherung (5 Min.):}
    Sicherung der Arbeitsartefakte (Branch-Stand, Test-/Lint-Ausgaben und Chat-Export) als Grundlage der späteren Artefaktanalyse.
    
    \item  \textbf{Debriefing (5 Min.):}
    Kurzes Abschlussgespräch, um offene Fragen zu klären und qualitatives Feedback einzuholen, das im Survey nicht abgebildet werden konnte.
\end{enumerate}

\section{Operationalisierung der Metriken}

Um die Hypothesen zu prüfen, wurden die abstrakten Beobachtungsdimensionen (Aufgabenfortschritt/Validierung sowie Interaktion) in nachvollziehbare, artefaktgebundene Indikatoren übersetzt.

\subsection{Aufgabenfortschritt und Validierung (Pipeline-Status)}
Die Ergebnislogik der Arbeit basiert auf einer artefaktgestützten Pipeline-Klassifikation pro Task und Branch. Zur Trennung von Intention, Implementationsartefakt und testgebundener Bestätigung werden die Statuskategorien  \textit{Attempted},  \textit{Implemented\textsubscript{static}} und  \textit{Verified} verwendet (Definitionen und Trichter-Evidenzstufe  \textit{Implemented$^\ast$} in Abschnitt~\ref{sec:results_pipeline}).
Wichtig ist dabei:  \textit{Verified} ist strikt testgebunden (\enquote{besteht die aktuellen Tests}) und somit ein methodengebundenes Diagnose-Signal, kein allgemeines Qualitätsurteil.

\subsection{Code- und Test-Artefakte (Surrogatindikatoren)}
Zur Einordnung der Umsetzungsstände werden zusätzlich struktur- und testnahe Artefakte herangezogen. Dazu zählen u.a. heuristische Strukturindikatoren und aggregierte Testausgänge aus den Auswertungsartefakten (z.\,B. \texttt{evaluation/deep\_code\_analysis.json}, \texttt{evaluation/code\_quality\_metrics.json}). Diese Größen werden durchgängig als Surrogate berichtet und nicht als vollständige Qualitätsmessung verstanden.

\subsection{Chat-Artefakte (Interaktion im Vibe-Coding)}
Die Interaktion wird über Chat-Exports und daraus abgeleitete deskriptive Indikatoren beschrieben (z.\,B. Wortanzahl, Prompt-Antwort-Zyklen, Fehlermarker-Iterationen). Die Auswertung bleibt auf Musterbeschreibung und Triangulation mit Pipeline- und Code/Test-Artefakten beschränkt; Chat-Intensität wird nicht als Effizienz- oder Qualitätsmaß interpretiert.

\section{Güte der Untersuchung (Threats to Validity)}

Die Validität der Studie wird anhand der Kategorien von \textcite{cook1979quasi} diskutiert.

\subsection{Konstruktvalidität (Construct Validity)}
Messen wir wirklich das, was wir messen wollen?
\begin{itemize}
    \item  \textbf{Problem:} \enquote{Kompetenz} ist schwer zu quantifizieren.
    \item  \textbf{Mitigation:} Kompetenz wird als Analyseperspektive über eine auditierbare, diskrete Pre-Survey-Heuristik operationalisiert (Tabelle~\ref{tab:competence-rubric}). Es wird explizit keine psychometrische Validierung beansprucht; die Einordnung wird über Triangulation mit Artefaktspuren (Pipeline/Code/Test/Chat) abgesichert.
\end{itemize}

\subsection{Interne Validität (Internal Validity)}
Können die beobachteten Effekte eindeutig der unabhängigen Variable (Kompetenz) zugeschrieben werden?
\begin{itemize}
    \item  \textbf{Selection Bias:} Die Zuteilung zu den Gruppen basierte auf Selbsteinschätzung. Um dies zu mitigieren, wurde der Algorithmus zur Klassifizierung strikt angewendet.
    \item  \textbf{Experimenter Bias (Rosenthal-Effekt):} Die Erwartungshaltung des Versuchsleiters könnte die Teilnehmenden beeinflussen.  \textbf{Gegenmaßnahme:} Standardisierte Instruktionen und minimale Interaktion während der Coding-Phase.
    \item  \textbf{Compensatory Rivalry (John Henry Effect):} Wenn eine Gruppe sich benachteiligt fühlt, könnte sie sich mehr anstrengen. Da alle Gruppen Copilot nutzten, ist dieses Risiko gering. Dennoch könnten Teilnehmende versuchen, besonders \enquote{sorgfältig} oder \enquote{kontrolliert} zu arbeiten (z.\,B. durch mehr manuelles Prüfen), was den beobachteten Fortschritt unter Zeitlimit beeinflussen kann \parencite{saretsky1972johnhenry}.
\end{itemize}

\subsection{Externe Validität (External Validity)}
Sind die Ergebnisse verallgemeinerbar?
\begin{itemize}
    \item  \textbf{Hawthorne-Effekt:} Das Wissen, beobachtet zu werden, verändert das Verhalten. Dies ist in Laborstudien typischerweise nicht vollständig vermeidbar, wurde aber durch die realistische Aufgabenstellung (Flow-Erleben) abgemildert.
    \item  \textbf{Ecological Validity:} Die Nutzung einer Brownfield-Umgebung erhöht die Übertragbarkeit auf Unternehmenskontexte im Vergleich zu synthetischen LeetCode-Aufgaben deutlich.
\end{itemize}

\subsection{Reliabilität (Reliability)}
Sind die Ergebnisse reproduzierbar?
\begin{itemize}
    \item  \textbf{Reproduzierbarkeit der Auswertung:}\\
    Die Pipeline-Status und Artefakt-Surrogate werden aus versionierten Branches und gespeicherten Auswertungsartefakten abgeleitet.
    Die Klassifikation folgt festen, dokumentierten Regeln (Attempted, Implemented\textsubscript{static} und Verified; Abschnitt~\ref{sec:results_pipeline}) und ist damit prinzipiell wiederholbar.
\end{itemize}

\section{Datenanalyse-Strategie}

Die Auswertung folgt einem artefaktbasierten, überwiegend deskriptiven Mixed-Methods-Ansatz:

\begin{enumerate}
    \item  \textbf{Deskriptive Quantifizierung:} Pipeline-Status (Attempted, Implemented\textsubscript{static} und Verified; sowie Implemented$^\ast$ für die Trichterdarstellung) werden pro Task und Branch berichtet. Ergänzend werden aggregierte Surrogatindikatoren aus Code- und Test-Artefakten ausgewiesen.
    \item  \textbf{Musteranalyse \& Triangulation:} Chat- und Strukturmuster werden als qualitative Befundlinien beschrieben und mit Pipeline- und Artefaktspuren trianguliert, um Integrations- und Validierungsengpässe im Brownfield-Kontext einzuordnen.
\end{enumerate}

Dieser methodische Aufbau stellt sicher, dass nicht nur das \emph{Ergebnis} (Code), sondern auch der \emph{Prozess} (Interaktion) beleuchtet wird.

\chapter{Technische Realisierung der Testumgebung}
\label{chap:implementierung}

Dieses Kapitel dokumentiert die technische Implementierung der \emph{BuyIn Todo Sandbox}, die als experimentelle Plattform für diese Studie dient. Es wird erläutert, wie die Brownfield-Umgebung architektonisch aufgebaut ist, um realistische Unternehmensanforderungen zu simulieren, und welche technischen Maßnahmen ergriffen wurden, um eine reproduzierbare Datenerhebung zu gewährleisten.

\section{Anforderungen an die Sandbox}

Um valide Aussagen über Vibe-Coding in Unternehmenskontexten treffen zu können, musste die Testumgebung folgende Anforderungen erfüllen:

\begin{enumerate}
    \item \textbf{Realismus (Brownfield):} Die Codebasis darf nicht leer sein (\enquote{Greenfield}), sondern muss bestehende Strukturen, technische Schulden und Konventionen enthalten, die der Agent und der Mensch verstehen müssen.
    \item \textbf{Komplexität:} Der Tech-Stack muss modern, aber komplex genug sein, um triviale Lösungen zu verhindern.
    \item \textbf{Isolierbarkeit:} Die Umgebung muss für jeden Teilnehmenden identisch zurückgesetzt werden können (Containerisierung).
    \item \textbf{Messbarkeit:} Das System muss automatisierte Tests bereitstellen, um den funktionalen Fortschritt objektiv zu messen.
\end{enumerate}

\section{Architektur und Tech-Stack: Kognitive Dimensionen}

Die Sandbox ist als klassische Full-Stack-Webanwendung konzipiert. Die Architekturwahl hat direkte Auswirkungen auf die kognitive Belastung der Teilnehmenden, analysiert durch das Framework der \emph{Cognitive Dimensions of Notation} \parencite{green1989cognitive}.

\subsection{Backend: Node.js \& Express (Layered Architecture)}
Das Backend basiert auf Node.js mit dem Express-Framework und folgt einer strikten Schichtenarchitektur:
\begin{itemize}
    \item \textbf{Routes:} Definition der API-Endpunkte.
    \item \textbf{Controllers:} Verarbeitung der HTTP-Requests und Validierung.
    \item \textbf{Services:} Kapselung der Geschäftslogik.
    \item \textbf{Repositories:} Datenzugriffsschicht.
\end{itemize}
\textbf{Kognitive Dimension - Viscosity (Zähigkeit):} Diese Trennung erhöht die \emph{Viscosity} für Novizen, da eine Änderung an mehreren Stellen (Datei-Hopping) durchgeführt werden muss. Für Copilot hingegen ist diese Struktur vorteilhaft, da sie klare semantische Grenzen zieht, die das Modell über den Dateinamen und Pfadkontext erkennen kann.

\subsection{Frontend: React \& Vite}
Das Frontend nutzt React mit TypeScript. Die Komponenten sind funktional aufgebaut und nutzen Hooks.
\textbf{Kognitive Dimension - Hidden Dependencies:} Die Verwendung von \emph{Custom Hooks} und \emph{Context API} erzeugt \emph{Hidden Dependencies}. Ein Low-Coder sieht im UI-Code nicht sofort, woher die Daten kommen. Dies testet die Fähigkeit des Agenten, diese Abhängigkeiten aufzulösen, und die Fähigkeit des Nutzers, dem Agenten den richtigen Kontext zu geben.

\section{Technische Funktionsweise von Copilot in der Sandbox}

Um zu verstehen, warum die Sandbox Copilot fordert, muss die technische Funktionsweise betrachtet werden.

\subsection{Prompt Construction und Context Retrieval}
Copilot sendet nicht nur den Code vor dem Cursor an das Modell. Der \emph{Copilot Proxy} in der IDE sammelt Kontext aus:
\begin{itemize}
    \item \textbf{Jaccard Similarity:} Offene Tabs werden nach Ähnlichkeit zum aktuellen Code gescannt.
    \item \textbf{Semantic Search:} Vektorisierte Suche im Projekt (bei neueren Versionen).
\end{itemize}
In der Sandbox sind die Aufgaben so gestaltet, dass die Lösung oft in einer \emph{anderen} Datei liegt (z.B. muss der Controller angepasst werden, um den Service zu nutzen). Wenn der Nutzer diese Datei nicht öffnet, fehlt dem Modell der Kontext. Dies misst indirekt die Kompetenz des Nutzers, dem Agenten die richtigen „Hinweise“ zu geben.

\subsection{Token Prediction und Ranking}
Das Modell generiert mehrere Vorschläge (Beams) und rankt diese nach Wahrscheinlichkeit. In einer Brownfield-Umgebung ist die Wahrscheinlichkeit für „korrekten“ Code höher, wenn er dem Stil der Umgebung entspricht. Ein High-Coder erkennt, ob der Vorschlag den Stil bricht (z.B. \texttt{var} statt \texttt{const}), ein Low-Coder nicht.

\section{Design der Aufgaben (Workload Grading)}

Die Aufgaben sind so gestaffelt, dass sie die kognitive Last sukzessive erhöhen:

\begin{table}[h]
    \centering
    \begin{tabular}{|l|l|l|}
    \hline
    \textbf{Task} & \textbf{Beschreibung} & \textbf{Kognitive Anforderung} \\ \hline
    1 & Persistenz (JSON) & Verstehen des Repository-Patterns (Backend). \\ \hline
    2 & User-Accounts & Einführung neuer Entitäten, Anpassung mehrerer Layer. \\ \hline
    3 & Kategorien (UI) & Frontend-Logik, State-Management. \\ \hline
    4 & Attachments & Umgang mit Binärdaten, API-Erweiterung. \\ \hline
    5 & Kalender & Komplexe Logik, Datumsberechnung (hoher Intrinsic Load). \\ \hline
    6 & Email-Verifikation & Integration externer Konzepte (Mocking). \\ \hline
    7 & Performance & Optimierung, kein funktionales Feature (Experten-Task). \\ \hline
    8 & ICS-Export & Standard-Format, reines Mapping (Erholung). \\ \hline
    \end{tabular}
    \caption{Aufgaben-Design und kognitive Anforderungen}
    \label{tab:task-design}
\end{table}

\section{Automatisierung der Evaluation}

Um die Auswertung der Ergebnisse zu standardisieren, wurde ein Evaluations-Framework entwickelt.

\subsection{Das Analyse-Skript (\texttt{run.js})}
Ein Node.js-Skript automatisiert die Prüfung der Lösungen. Es führt folgende Schritte aus:
\begin{enumerate}
    \item \textbf{Test-Execution:} Ausführung der Unit-Tests (Vitest) und End-to-End-Tests (Playwright).
    \item \textbf{Linting:} Prüfung auf Code-Style-Verletzungen mittels ESLint.
    \item \textbf{Report-Generierung:} Zusammenfassung der Ergebnisse in einer JSON-Datei.
\end{enumerate}

\subsection{Grenzen der LLM-gestützten Code-Analyse}
Zusätzlich zur rein funktionalen Prüfung wird ein LLM-Prompt-Template verwendet, um qualitative Aspekte des Codes zu bewerten. Hierbei muss kritisch angemerkt werden, dass LLMs zu \emph{Halluzinationen} neigen können, also Fehler finden, die nicht existieren, oder echte Probleme übersehen. Daher dient dieser Score nur als Indikator und wird stichprobenartig manuell validiert.

\section{Relevanz für den Unternehmenskontext (BuyIn)}

Die gewählte Architektur und Methodik spiegelt reale Herausforderungen bei BuyIn wider:
\begin{itemize}
    \item \textbf{Legacy-Integration:} Wie bei BuyIn müssen Entwickler oft in gewachsenen Strukturen arbeiten, die nicht perfekt dokumentiert sind. Vibe-Coding kann hier helfen, den „impliziten Kontext“ schneller zu erschließen.
    \item \textbf{Security \& Compliance:} Der Einsatz von Copilot in Enterprise-Umgebungen birgt Risiken (Datenabfluss, Lizenzverletzungen). Die Sandbox simuliert dies durch Aufgaben, bei denen sensible Daten (User-Accounts) behandelt werden müssen.
    \item \textbf{Skalierbarkeit:} Die modulare Struktur der Sandbox zeigt, wie Copilot in größeren Teams (wo Modularisierung Pflicht ist) performt, im Gegensatz zu kleinen Skript-Projekten.
\end{itemize}

\section{Containerisierung und Reproduzierbarkeit}

Die gesamte Umgebung ist mittels Docker containerisiert. Dies stellt sicher, dass alle Teilnehmenden – unabhängig von ihrem Betriebssystem – unter identischen Bedingungen arbeiten und technische Störfaktoren (z.B. unterschiedliche Node-Versionen) eliminiert werden.

\chapter{Ergebnisse}
\label{chap:results}

Zur Orientierung gilt: Die Pipeline-Status \textit{Attempted}, \textit{Implemented\textsubscript{static}} und \textit{Verified} beschreiben unterschiedliche Evidenzebenen desselben Arbeitsprozesses und sind keine Qualitäts- oder Kompetenzbewertung. Insbesondere gilt: \textit{Verified} ist nicht gleich \enquote{besserer Entwickler}. \textit{Implemented\textsubscript{static}} meint strukturelle Artefakte in der Codebasis; \textit{Verified} ist strikt testgebunden.

\section{Datengrundlage und Toolkontext}
\label{sec:results_context}

Dieses Unterkapitel beschreibt die Datengrundlage der Untersuchung sowie den technischen Kontext der durchgeführten \textit{VibeCoding-Challenge}. Die Analyse basiert auf einem Datensatz von $N=18$ Probanden, die im Rahmen einer kontrollierten Coding-Challenge Aufgaben zur Entwicklung einer Todo-Applikation bearbeiteten.

\subsection{Studiensetting und Rahmenbedingungen}
Die Studie wurde unter einheitlichen Rahmenbedingungen durchgeführt, um die Vergleichbarkeit der Ergebnisse zu gewährleisten:
\begin{itemize}
    \item \textbf{Zeitlimit:} Alle Probanden hatten ein striktes Zeitlimit von 60 Minuten für die Bearbeitung der Aufgaben.
    \item \textbf{Tooling:} Als Entwicklungsumgebung diente Visual Studio Code mit der GitHub Copilot Extension.
    \item \textbf{KI-Modell:} Als zugrundeliegendes Large Language Model (LLM) wurde \textbf{Claude 3.5 Sonnet} verwendet. Die Interaktion erfolgte über das Chat-Interface von GitHub Copilot, wobei das Modell explizit als Backend konfiguriert war.
    \item \textbf{Aufgabenstellung:} Die Probanden erhielten eine Liste von 8 Aufgaben (T1--T8) mit steigender Komplexität, beginnend bei Backend-Persistenz bis hin zu komplexen Frontend-Features und Export-Funktionen.
\end{itemize}

\subsection{Datenquellen und Artefakte}
Die empirische Auswertung stützt sich auf drei primäre Datenquellen, die eine Triangulation der Ergebnisse ermöglichen:

\begin{enumerate}
    \item \textbf{Code-Artefakte (Git-Repositories):}
    Die finalen Implementierungen der Probanden liegen als separate Git-Branches vor. Diese Branches enthalten den vollständigen Quellcode zum Zeitpunkt des Zeitablaufs.
    \begin{itemize}
        \item Pfad: Remote-Branches mit dem Suffix \texttt{*-vibe} (z.B. \texttt{origin/august-martin-vibe}).
    \end{itemize}

    \item \textbf{Chat-Logs:}
    Die vollständigen Interaktionsprotokolle zwischen Proband und KI-Assistent wurden exportiert. Diese Protokolle enthalten alle Prompts der Nutzer sowie die Antworten und Code-Generierungen der KI.
    \begin{itemize}
        \item Pfad: \texttt{thesis/chats/} (Format: \texttt{.docx}).
    \end{itemize}

    \item \textbf{Survey-Daten:}
    Im Anschluss an die Challenge füllten die Probanden einen Fragebogen zur Selbsteinschätzung und subjektiven Wahrnehmung der KI-Interaktion aus.
    \begin{itemize}
        \item Pfad: \texttt{thesis/surveys/Probanden-Einstufungsformular (Responses).xlsx}.
    \end{itemize}
\end{enumerate}

\subsection{Dateninventar und Vollständigkeit}
Tabelle \ref{tab:data_inventory} zeigt das Inventar der verfügbaren Datenartefakte. Insgesamt liegen für 18 Probanden Daten vor.

\begin{table}[h]
    \centering
    \begin{tabular}{|l|c|l|}
    \hline
    \textbf{Artefakt-Typ} & \textbf{Anzahl} & \textbf{Anmerkung} \\
    \hline
    Git-Branches (\texttt{*-vibe}) & 16 & 2 Probanden arbeiteten lokal/ohne Push* \\
    Chat-Logs & 18 & Vollständig für alle Teilnehmer \\
    Survey-Einträge & 18 & Vollständig ($N=18$) \\
    \hline
    \end{tabular}
    \caption{Inventar der Datenartefakte. (*Für die fehlenden Remote-Branches liegen lokale Stände oder Chat-Rekonstruktionen vor, die für die Analyse genutzt wurden.)}
    \label{tab:data_inventory}
\end{table}

\subsection{Zuordnung und Zusammenführung der Daten}
Die Zusammenführung der Datenquellen erfolgte über eine Mapping-Tabelle, die die Klarnamen aus dem Survey und den Chat-Dateinamen den entsprechenden Git-Branches zuordnet.
\begin{itemize}
    \item \textbf{Survey $\leftrightarrow$ Branch:} Das Survey-Feld "Branch-Name" sowie der Name des Probanden dienten als Schlüssel.
    \item \textbf{Chat $\leftrightarrow$ Branch:} Die Dateinamen der Chat-Logs (z.B. \texttt{Max\_Helling.docx}) wurden normalisiert und den entsprechenden Branches (z.B. \texttt{helling-max-vibe}) zugeordnet.
\end{itemize}
Die Survey-Daten enthalten zudem eine Einteilung in Kompetenzgruppen (No-Code, Low-Code, Mid-Code, High-Code), die in der späteren Analyse als vergleichendes Merkmal herangezogen wird.

\subsection{Reproduzierbarkeit und Auswertungswerkzeuge}
Um die Ergebnisse objektiv und reproduzierbar zu gestalten, kommen automatisierte Analysewerkzeuge zum Einsatz, deren Resultate im Ordner \texttt{evaluation/} persistiert werden:

\begin{itemize}
    \item \textbf{Statische Code-Analyse:} Einsatz von \texttt{ESLint} zur Überprüfung der Code-Qualität und Einhaltung von Standards.
    \item \textbf{Automatisierte Tests:} Ausführung der vorhandenen Test-Suite (basierend auf \texttt{Jest/Vitest}), um die funktionale Korrektheit der Implementierungen (Task Completion) zu verifizieren.
    \item \textbf{Chat-Parsing:} Algorithmische Auswertung der Chat-Logs zur Bestimmung von "Attempted"-Tasks und Identifikation von Prompting-Mustern.
\end{itemize}

\subsection{Abgrenzung zu Kapitel 7}
Das vorliegende Kapitel 6 beschränkt sich auf die deskriptive Darstellung der Ergebnisse und die objektive Auswertung der Daten. Eine Interpretation dieser Ergebnisse, die Diskussion von Ursache-Wirkungs-Beziehungen sowie die Beantwortung der Forschungsfragen erfolgen im anschließenden Kapitel 7 (Diskussion).

\section{Task-Pipeline-Analyse}
\label{sec:results_pipeline}

In diesem Abschnitt wird der Fortschritt der Probanden entlang der Aufgaben-Pipeline (T1--T5) dargestellt. Berichtete werden Häufigkeiten für Intention (Chat), Umsetzung (Code-Artefakte) und den Statusausweis \textit{Execution observed} (Integrationsläufe).

\subsection{Methodische Definitionen}
Für die Auswertung werden drei Statuskategorien verwendet:
\begin{itemize}
    \item  \textbf{Attempted (Versucht):} Ein Task gilt als  \textit{Attempted}, wenn er im Chat-Verlauf explizit als Ziel, Anforderung oder nächster Arbeitsschritt formuliert wurde. Die bloße Existenz von Dateien ohne entsprechenden Kontext im Chat genügt nicht.
    \item  \textbf{Implemented\textsubscript{static} (Umgesetzt, Artefakt/Heuristik):}\\
    Ein Task gilt als  \textit{Implemented\textsubscript{static}}, wenn er in den Auswertungsartefakten als implementiert markiert ist (heuristische, statische Indikatoren; z.\,B.\ Feld \texttt{implemented}\\
    in \path{evaluation/task_pipeline_v3.json}\\
    bzw.\ Implementationsklassen in \mbox{\path{evaluation/deep_code_analysis.json}}).
    \item  \textbf{Execution observed (Statusausweis):} Ein Task gilt als  \textit{Execution observed}, wenn die zugehörigen automatisierten Integrationsläufe erfolgreich durchlaufen wurden (Feld \texttt{verified} in \path{evaluation/task_pipeline_v3.json}).
\end{itemize}

\noindent \textbf{Abgeleitete Evidenzstufe.} In Kapitel~6 wird eine abgeleitete Evidenzstufe\\
\BeginAccSupp{unicode,method=pdfstringdef,ActualText={extitImplemented* verwendet:}}
\csname textit\endcsname{Implemented$^\ast$} verwendet:
\EndAccSupp{}
\[
\begin{aligned}
    \csname textit\endcsname{Implemented$^\ast$} &:=  \textit{Attempted} \wedge (\textit{Implemented\textsubscript{static}} \vee  \textit{Execution observed}).
\end{aligned}
\]
Damit gilt per Konstruktion:  \textit{Execution observed} $\subseteq$  \textit{Implemented$^\ast$}.  \textit{Execution observed} ist jedoch nicht notwendigerweise ein Subset von  \textit{Implemented\textsubscript{static}}:  \textit{Implemented\textsubscript{static}} basiert auf heuristischen, strukturellen Erkennungsregeln; unter den gegebenen Rahmenbedingungen ist es möglich, dass funktionsfähige Implementationen nicht als  \textit{Implemented\textsubscript{static}} erkannt werden (z.\,B. abweichende Benennungen, alternative Architekturpfade oder Implementationen außerhalb der erwarteten Muster).  \textit{Execution observed} ist ein methodengebundenes Diagnose-Signal;  \textit{Implemented\textsubscript{static}} ist ein Artefakt-Signal.

\noindent \textbf{Hinweis zur Ausführungsmethodik.} Für die Pipeline-Auswertung wurden die automatisierten Integrationsläufe (insbesondere T2, sowie in Teilen T3--T5) als API-Integrationsläufe (\textit{Supertest} gegen Express) gestaltet. \textit{Execution observed} bedeutet in dieser Studie explizit: „durchläuft die aktuellen Integrationsläufe“ \textendash{} nicht „fachlich vollständig korrekt gelöst“.

\subsection{Pipeline-Ergebnisse (T1--T5)}
\begin{table}[ht]
\centering
\begingroup
\newcommand{\NoExtract}[1]{\BeginAccSupp{method=escape,ActualText={}}#1\EndAccSupp{}}
\newcommand{\ExtractOnly}[1]{\smash{\rlap{\BeginAccSupp{unicode,method=pdfstringdef,ActualText={#1}}\textcolor{white}{.}\EndAccSupp{}}}}
\def\TaskPipelineTableActualText{Task^^JAttempted (𝑛) Implemented* (𝑛) Execution observed (𝑛)^^JT1: Persistenz^^J18 (100\%)^^J17 (94\%)^^J10 (56\%)^^JT2: Auth^^J18 (100\%)^^J17 (94\%)^^J4 (22\%)^^JT3: Kategorien^^J15 (83\%)^^J13 (72\%)^^J13 (72\%)^^JT4: Attachments^^J13 (72\%)^^J8 (44\%)^^J0 (0\%)^^JT5: Kalender^^J13 (72\%)^^J0 (0\%)^^J0 (0\%)^^J^^JTabelle 6.2: Task-Pipeline: Von der Intention zur testgebundenen Bestätigung (𝑁 = 18,^^J𝑛 = Anzahl Branches).^^JextitAttempted stammt aus Chat-Signalen. Implemented* ist definiert als: Attempted und^^J(Implementedstatic oder Execution observed). Execution observed stammt aus V3-Integrationsläufen^^J(V3 bezeichnet die für diese Studie definierte Integrationslauf-Konfiguration).}
\NoExtract{\begin{tabular}{|l|c|c|c|}
\hline
    \csname textbf\endcsname{Task} &  \textbf{Attempted ($n$)} &  \textbf{Implemented$^\ast$ ($n$)} &  \textbf{Execution observed ($n$)} \\
\hline
T1: Persistenz & 18 (100\%) & 17 (94\%) & 10 (56\%) \\
T2: Auth & 18 (100\%) & 17 (94\%) & 4 (22\%) \\
T3: Kategorien & 15 (83\%) & 13 (72\%) & 13 (72\%) \\
T4: Attachments & 13 (72\%) & 8 (44\%) & 0 (0\%) \\
T5: Kalender & 13 (72\%) & 0 (0\%) & 0 (0\%) \\
\hline
\end{tabular}}
\ExtractOnly{\TaskPipelineTableActualText}
\NoExtract{\caption{Task-Pipeline: Von der Intention zur testgebundenen Bestätigung ($N = 18$, $n$ = Anzahl Branches). \textit{Attempted} stammt aus Chat-Signalen. \textit{Implemented$^\ast$} ist definiert als: \textit{Attempted} und (\textit{Implemented\textsubscript{static}} oder \textit{Execution observed}). \textit{Execution observed} stammt aus V3-Integrationsläufen (V3 bezeichnet die für diese Studie definierte Integrationslauf-Konfiguration).}}
\label{tab:task_pipeline_v2}
\endgroup
\end{table}

\noindent \textbf{Messung.} Berichtete werden je Task (T1--T5) die Häufigkeiten der Pipeline-Status \textit{Attempted} (Chat-Signal), \textit{Implemented$^\ast$} (kombinierte Evidenzstufe) und \textit{Execution observed} (methodengebundenes Signal) über $N=18$ Branches.
\noindent \textbf{Hauptbefund.} Für T1/T2 liegt \textit{Attempted} jeweils in 18/18 Branches vor, während \textit{Execution observed} für T4/T5 in 0/18 Branches erreicht wird.
\noindent \textbf{RQ-Bezug.} Die Statushäufigkeiten bilden die deskriptive Ergebnisbasis für RQ1 und werden in Kapitel~7 für die Einordnung der Befunde und die Pipeline-Diskussion verwendet.

\begin{figure}[ht]
\centering
\makebox[\linewidth][c]{%
\begin{tikzpicture}
\begin{axis}[
    xbar,
    width=0.92\linewidth,
    height=7.2cm,
    xmin=0,
    xmax=18,
    xtick={0,2,...,18},
    scaled x ticks=false,
    tick label style={/pgf/number format/fixed,/pgf/number format/precision=0},
    xticklabel style={font=\scriptsize},
    xlabel={Anzahl Branches ($N=18$)},
    xlabel style={font=\small},
    symbolic y coords={T1: Persistenz,T2: Auth,T3: Kategorien,T4: Attachments,T5: Kalender},
    ytick=data,
    y dir=reverse,
    bar width=6pt,
    enlarge y limits=0.45,
    trim axis left,
    trim axis right,
    axis lines=left,
    xmajorgrids,
    grid style={black!8, line width=0.15pt},
    tick style={black!50},
    legend style={
        at={(0.5,-0.18)},
        anchor=north,
        legend columns=3,
        font=\small,
        draw=none
    },
    point meta=explicit symbolic,
    nodes near coords,
    nodes near coords style={
        font=\scriptsize,
        anchor=west,
        xshift=3pt,
        text=black!80
    },
    clip=false,
]
    % Attempted (hellgrau)
    \addplot[fill=black!15, draw=black!55, line width=0.2pt, bar shift=9pt]
    coordinates {(18,T1: Persistenz) [18] (18,T2: Auth) [18] (15,T3: Kategorien) [15] (13,T4: Attachments) [13] (13,T5: Kalender) [13]};

    % Implemented* (mittelgrau, visuell dominanter)
    \addplot[fill=black!55, draw=black, line width=0.35pt, bar width=7pt, bar shift=0pt]
    coordinates {(17,T1: Persistenz) [17] (17,T2: Auth) [17] (13,T3: Kategorien) [13] (8,T4: Attachments) [8] (0,T5: Kalender) [{}]};

    % Execution observed (dunkelgrau/schwarz; zusätzlich Schraffur)
    \addplot[fill=black!85, draw=black, line width=0.4pt, pattern=north east lines, pattern color=black, bar shift=-9pt]
    coordinates {(10,T1: Persistenz) [10] (4,T2: Auth) [4] (13,T3: Kategorien) [13] (0,T4: Attachments) [{}] (0,T5: Kalender) [{}]};

    % 0-Werte (ohne Textlabel) optional als Marker am Achsenstart
    \addplot[only marks, mark=o, mark size=2.6pt, draw=black, mark layer=foreground, forget plot, bar shift=0pt]
    coordinates {(0,T5: Kalender)};
    \addplot[only marks, mark=x, mark size=3.0pt, draw=black, mark layer=foreground, forget plot, bar shift=-9pt]
    coordinates {(0,T4: Attachments) (0,T5: Kalender)};

    \legend{Attempted, Implemented$^\ast$, Execution observed}
\end{axis}
\end{tikzpicture}
}%
\caption{Pipeline-Ergebnisse für T1--T5 ($N=18$) als gruppierte Balken:  \textit{Attempted} (Chat-Signal),  \textit{Implemented$^\ast$} (definiert als:  \textit{Attempted} und (\textit{Implemented\textsubscript{static}} oder \textit{Execution observed})) und  \textit{Execution observed} (testgebunden: bestandene V3-Integrationsläufe). V3 bezeichnet die für diese Studie definierte Integrationslauf-Konfiguration; die Werte entsprechen Tabelle~\ref{tab:task_pipeline_v2}.}
\label{fig:task_pipeline_funnel_t1_t5}
\end{figure}

\noindent \textbf{Messung.} Visualisiert werden die in Tabelle~\ref{tab:task_pipeline_v2} berichteten Statushäufigkeiten (\textit{Attempted}, \textit{Implemented$^\ast$}, \textit{Execution observed}) als Balken pro Task.
\noindent \textbf{Hauptbefund.} Über T1--T5 liegen \textit{Implemented$^\ast$} und \textit{Execution observed} unter \textit{Attempted} und erreichen bei T4/T5 den Wert 0.
\noindent \textbf{RQ-Bezug.} Die Abbildung ist eine komprimierte Darstellung der Pipeline-Ergebnisse und ergänzt die Ergebniszusammenfassung zu RQ1 in Kapitel~7.

\subsection{Kurzbefunde pro Task (T1--T5)}
\noindent \textbf{T1 (Persistenz).} T1 ist in 18/18 Branches  \textit{Attempted}, in 17/18 Branches  \textit{Implemented$^\ast$} und in 10/18 Branches  \textit{Execution observed}.

\noindent \textbf{T2 (User Accounts/Auth).} T2 ist in 18/18 Branches  \textit{Attempted}, in 17/18 Branches  \textit{Implemented$^\ast$} und in 4/18 Branches  \textit{Execution observed}.

\noindent \textbf{T3 (Kategorien).} T3 ist in 15/18 Branches  \textit{Attempted}, in 13/18 Branches  \textit{Implemented$^\ast$} und in 13/18 Branches  \textit{Execution observed}. In \path{evaluation/task_pipeline_v3.json} ist T3 dabei in 7/18 Branches als \texttt{implemented} (\textit{Implemented\textsubscript{static}}) markiert.

\noindent \textbf{T4 (Attachments).} T4 ist in 13/18 Branches  \textit{Attempted}, in 8/18 Branches  \textit{Implemented$^\ast$} und in 0/18 Branches  \textit{Execution observed}.

\noindent \textbf{T5 (Kalender/Multi-Day).} T5 ist in 13/18 Branches  \textit{Attempted}, in 0/18 Branches  \textit{Implemented$^\ast$} und in 0/18 Branches  \textit{Execution observed}.

% DATENQUELLEN:
% Chat-Logs (für "Attempted"-Status)
% Code-Analyse (für "Implemented"-Status)
% Test-Reports (für "Execution observed"-Status)
%
% HINWEISE:
% Definition "Attempted": Explizite Erwähnung im Chat
% Definition "Implemented": Code vorhanden
% Definition "Execution observed": Tests bestanden
% Fokus auf die Statushäufigkeiten über die Tasks hinweg

\section{Detaillierte Lösungsmuster pro Task (T1--T5)}
\label{sec:results_strategies}

\BeginAccSupp{unicode,method=pdfstringdef,ActualText={Dieses Unterkapitel beschreibt typische Lösungsmuster, die in den 𝑁 = 18 *-vibeBranches für die Tasks T1–T5 beobachtet wurden. Grundlage sind Repository-Struktur^^Jund Implementationsartefakte (Backend/Frontend, Datenhaltung), Chat-basierte Attempted-Signale sowie Survey-Angaben (aggregiert, ohne Klarnamen). Die Darstellung ist^^Jrein deskriptiv.}}
Dieses Unterkapitel beschreibt typische Lösungsmuster, die in den $N = 18$ *-vibeBranches für die Tasks T1--T5 beobachtet wurden. Grundlage sind Repository-Struktur und Implementationsartefakte (Backend/Frontend, Datenhaltung), Chat-basierte  \textit{Attempted}-Signale sowie Survey-Angaben (aggregiert, ohne Klarnamen). Die Darstellung ist rein deskriptiv.
\EndAccSupp{}

\subsection*{Einordnung der Datenbasis (kurz)}
T1 und T2 wurden im Chat durchgängig adressiert (18/18), T3 in 15/18 sowie T4/T5 in jeweils 13/18 Fällen. Entsprechend nimmt die Abdeckung der Chat-basierten  \textit{Attempted}-Signale von T1/T2 zu T4/T5 ab. In den Surveys werden Zeitdruck und Debugging am häufigsten als Hürden genannt (je 5/18), gefolgt von Aufgabenverständnis (3/18). Die Selbsteinstufung verteilt sich über No-Code (5), Low-Code (4), Mid-Code (5) und High-Coder (4).

\subsection*{T1 -- Persistente Todos}
\begin{enumerate}
  \item  \textbf{Ziel des Tasks (kurz, technisch).} Persistente Speicherung von Todos über Prozess-Neustarts hinweg, inkl. stabilem CRUD-Verhalten (Erstellen, Lesen, Aktualisieren, Löschen) und konsistenter Identifikatoren.

  \item  \textbf{Typische Implementierungsvarianten (mit Häufigkeiten).} Die häufigste Variante war eine SQLite-basierte Persistenz (14/18). Daneben traten MongoDB (1/18) und PostgreSQL (1/18) sowie verbleibende Varianten wie In-Memory/sonstige Ablagen (2/18) auf. In der Code-Struktur wird häufig ein Repository-Pattern (z.B. \texttt{TodoRepository}) mit einer klaren Abgrenzung zwischen Controller/Route und Datenzugriff verwendet.

  \item  \textbf{Beobachtete Stabilitätsmuster (konsistent vs. inkonsistent).} Konsistent sind konsistente ID-Schemata und deterministisches CRUD. Inkonsistent sind v.a. uneinheitliche ID-Felder (\texttt{id} vs. \texttt{\_id}) und inkonsistente Update-Semantik.

  \item  \textbf{Beobachtete Implementationsmuster.} 
  \begin{sloppypar}Häufig treten Repository-/Controller-Boilerplates und generische CRUD-Endpunkte auf. Randfälle (Eingabeprüfung, Fehlercodes, ID-Normalisierung) sind in den Artefakten nicht in allen Branches durchgängig konsistent umgesetzt.\end{sloppypar}

  \item  \textbf{Beobachtete Teilumsetzungen / offene Punkte.} In einzelnen Branches bleiben Randfälle (Eingabeprüfung, Fehlercodes, ID-Normalisierung) sowie Konsistenzanforderungen (einheitliche ID-Felder, deterministische Update-Semantik) unvollständig.
\end{enumerate}

\subsection*{T2 -- Benutzerkonten / Authentifizierung}
\begin{enumerate}
  \item  \textbf{Ziel des Tasks (kurz, technisch).} Nutzerregistrierung und Login mit Zugriffskontrolle auf geschützte Endpunkte (z.B. \texttt{GET /api/todos}) sowie konsistenter Authentifizierungsweitergabe (Token oder Session).

  \item  \textbf{Typische Implementierungsvarianten (mit Häufigkeiten).} JWT ist die häufigste Variante (11/18), gefolgt von Session/Cookie (4/18) und Hybrid-Ansätzen (1/18). Selten sind Sonderfälle (z.B. Hashing ohne klaren Token-Flow: 1/18) bzw. fehlende oder unklare Signale (1/18). Häufige Strukturmuster sind zentrale Auth-Middleware oder Kontrollen direkt im Controller.

  \item  \textbf{Beobachtete Stabilitätsmuster (konsistent vs. inkonsistent).} Konsistent sind eindeutige Token-/Session-Flows und konsistente Route-Guards. Inkonsistent sind inkonsistente Header-/Cookie-Nutzung und divergierende Erwartungen zwischen Frontend und Backend.

  \item  \textbf{Beobachtete Implementationsmuster.}\\
  Boilerplate-Flows (Register, Login, Protect) treten häufig auf. In den Artefakten finden sich sowohl Varianten mit Rollen-/Claims-Strukturen als auch minimal gehaltene Flows (z.B. ohne durchgängige Eingabeprüfung/Expiry). Häufig treten Kombinationen aus JWT und \texttt{bcrypt} auf.

  \item  \textbf{Beobachtete Teilumsetzungen / offene Punkte.} In einzelnen Branches sind Header-/Cookie-Nutzung und Route-Guards inkonsistent oder zwischen Frontend und Backend nicht durchgängig abgestimmt.
\end{enumerate}

\subsection*{T3 -- Kategorien}
\begin{enumerate}
  \item  \textbf{Ziel des Tasks (kurz, technisch).} Erweiterung des Todo-Datenmodells um Kategorisierung und Bereitstellung über API und UI, sodass Todos kategoriebasiert erstellt, gespeichert und angezeigt werden können.

  \item  \textbf{Typische Implementierungsvarianten (mit Häufigkeiten).} Ein separates Kategorie-Konzept (Model/Tabelle/Collection mit Relation via \texttt{categoryId}) dominiert (14/18). Inline-Kategorien als String sind selten (1/18). In 3/18 Branches ist das Kategorie-Konzept in Struktursignalen nicht klar erkennbar.

  \item  \textbf{Beobachtete Stabilitätsmuster (konsistent vs. inkonsistent).} Konsistent sind konsistente Persistenz und stabile Referenzen (IDs/Relation). Inkonsistent sind rein UI-seitige Kategorien oder inkonsistente Relation/Serialisierung.

  \item  \textbf{Beobachtete Implementationsmuster.} Häufig treten zusätzliche \texttt{Category}-Modelle und Routen auf. In einzelnen Branches werden Kategorien und Tags terminologisch vermischt; dabei variiert die Bezeichnung, während die Datenrepräsentation nicht in allen Fällen konsistent vereinheitlicht ist.

  \item  \textbf{Beobachtete Teilumsetzungen / offene Punkte.} In einzelnen Branches bleiben Kategorien rein UI-seitig, oder Persistenz/Serialisierung/Relation ist nicht durchgängig konsistent umgesetzt.
\end{enumerate}

\subsection*{T4 -- Dateianhänge (Attachments)}
\begin{enumerate}
  \item  \textbf{Ziel des Tasks (kurz, technisch).} Upload und Zuordnung von Dateien zu Todos (typischerweise \texttt{multipart/form-data}), Speicherung (Dateisystem oder DB) und Rückgabe verwertbarer Metadaten über die API.

  \item  \textbf{Typische Implementierungsvarianten (mit Häufigkeiten).} In 8/18 Branches ist ein Multer-/Upload-Ansatz erkennbar. In 10/18 Branches sind Attachments nicht oder nur indirekt oder unklar umgesetzt. Wo Upload existiert, ist Storage häufig lokal (Uploads-Ordner) oder metadatenbasiert (Dateiname/URL).

  \item  \textbf{Beobachtete Stabilitätsmuster (konsistent vs. inkonsistent).} Konsistent sind klare Upload-Verträge (Feldnamen/Metadaten) und stabile Zuordnung zum Todo. Inkonsistent sind uneinheitliche Response-Formate und fehlende Metadaten-Persistenz.

  \item  \textbf{Beobachtete Implementationsmuster.}\\
  Multer-Templates mit Boilerplate-Routen sind häufig. Daneben finden sich sowohl Varianten mit separaten Attachment-Entitäten als auch Upload-Varianten ohne stabilen Metadaten-Rückkanal.

  \item  \textbf{Beobachtete Teilumsetzungen / offene Punkte.} Wo Upload existiert, sind Response-Formate, Feldnamen und Metadatenpersistenz nicht in allen Branches konsistent.
\end{enumerate}

\subsection*{T5 -- Kalender / Multi-Day}
\begin{enumerate}
  \item  \textbf{Ziel des Tasks (kurz, technisch).} Abbildung von Todos im Kalender, inklusive Unterstützung von Zeiträumen (Multi-Day), sowie konsistente Query-/Filter-Semantik (z.B. Overlap-Logik) zur effizienten Darstellung in Kalenderansichten.

  \item  \textbf{Typische Implementierungsvarianten (mit Häufigkeiten).} In den Branches sind Range-Felder als Struktursignale verbreitet (z.B. \texttt{dueEndDate} neben \texttt{dueDate}). Gleichzeitig ist T5 in den (V3-)Integrationsläufen in 0/18 Branches als \textit{Execution observed} ausgewiesen.

  \item  \textbf{Beobachtete Stabilitätsmuster (konsistent vs. inkonsistent).} In den Artefakten treten zusätzliche Range-Felder häufig auf. Eine durchgängig erkennbare Range-Logik (Ende > Start) sowie eine konsistente Query-/Filter-Semantik (Overlap-Logik) ist in den Artefakten nicht in allen Branches erkennbar.

  \item  \textbf{Beobachtete Implementationsmuster.} Häufig treten UI-/Modell-Boilerplates (zusätzliche Datumsfelder, Kalender-Komponenten) sowie Range-Utilities auf. In den Artefakten ist der API-Vertrag (Parameter/Overlap-Regel) dabei nicht in allen Branches konsistent erkennbar.

  \item  \textbf{Beobachtete Teilumsetzungen / offene Punkte.} In den Artefakten sind Range-Felder und UI-Komponenten teilweise vorhanden, ohne dass eine durchgängig konsistente, durch \textit{Execution observed} ausgewiesene End-to-End-Semantik für Overlap/Filter und API-Vertrag vorliegt.
\end{enumerate}

\section{Code-Qualität und Metriken}
\label{sec:results_quality}

Dieses Unterkapitel präsentiert die objektiven Messwerte zur Code-Qualität, basierend auf automatisierten Analysetools.

% DATENQUELLEN:
% - ESLint Reports (Linting Errors)
% - Test Reports (Pass/Fail Raten)
% - Lighthouse/Playwright (falls vorhanden)
%
% HINWEISE:
% - Rein quantitative Darstellung
% - Keine Bewertung der Ursachen (das folgt in Kap. 7)

\section{Chat-Interaktionsmuster}
\label{sec:results_chat}

Die Analyse der Interaktionen zwischen Probanden und dem KI-Modell wird in diesem Abschnitt dargestellt.

% DATENQUELLEN:
% - Chat-Logs (Prompts, Antworten)
%
% HINWEISE:
% - Prompt-Qualität (Struktur, Kontext)
% - Debugging-Loops (Ping-Pong-Effekte)
% - Nutzung von Claude Sonnet 4.5 Spezifika

\section{Triangulation der Datenquellen}
\label{sec:results_triangulation}

Hier werden die Ergebnisse aus Code-Analyse, Chat-Logs und Test-Reports zusammengeführt, um Zusammenhänge sichtbar zu machen.

% DATENQUELLEN:
% - Aggregierte Daten aus allen vorherigen Analysen
%
% HINWEISE:
% - Zusammenhang zwischen Chat-Verhalten und Code-Qualität
% - Zusammenhang zwischen Strategie (z.B. DB-Wahl) und Testerfolg
% - Rein deskriptive Korrelationen, keine kausale Interpretation

\section{Zusammenfassung der Ergebnisse}
\label{sec:results_summary}

Kapitel~6 dokumentiert die Ergebnisse der Evaluation und stellt beobachtete Ausprägungen und Verteilungen dar, ohne diese zu bewerten oder zu erklären.
Die Ergebnisdarstellung basiert auf mehreren Datenquellen (Pipeline-Status \,\texttt{attempted}/\texttt{implemented}/\texttt{verified}, Code- und Test-Artefakte, Chat-Exports sowie Survey-Merkmale) und bleibt durchgängig deskriptiv.
Im Folgenden werden die zentralen Befundlinien systematisch zusammengeführt.

\paragraph{Pipeline-Ergebnisse entlang der Tasks (T1--T5).}
Über die Tasks hinweg zeigt sich ein gemeinsamer Trichter-Effekt: Attempted-Markierungen sind in der Breite vorhanden, während die Anteile der als Implemented$^\ast$ (Trichter-Evidenzstufe) bzw. Verified ausgewiesenen Ergebnisse mit zunehmendem Integrations- und Abhängigkeitsgrad abnehmen.
Dabei ist die Ausprägung dieses Trichters nicht in allen Branches gleichförmig, sondern zeigt Spannweiten zwischen sehr kurzem Task-Fortschritt und einem deutlich weiter reichenden Umsetzungsstand.
In der Gegenüberstellung der Aufgaben treten dabei zwei Bündel hervor:
\begin{itemize}
\item \textbf{Backend-nahe Tasks (T1/T2).} T1 ist in den Branches in unterschiedlichen Persistenzvarianten umgesetzt, wobei eine Variante dominiert und mehrere Alternativen vereinzelt auftreten. T2 wird in der Mehrzahl der Branches als Implemented$^\ast$ ausgewiesen und ist damit ein häufig realisiertes Backend-Element.
\item \textbf{Integrations- und Frontend-nahe Tasks (T3--T5).} Bei T3 und T4 nehmen die Implemented\textsubscript{static}-Ausweise ab. Zugleich werden bei einzelnen Tasks Diskrepanzen zwischen Implemented\textsubscript{static} (heuristische Artefakt-Markierung) und Verified sichtbar (ohne dass daraus in Kapitel~6 Schlussfolgerungen abgeleitet werden). T5 bleibt in den Artefakten durchgängig ohne Implemented\textsubscript{static}- und ohne Verified-Ausweise.
\end{itemize}
Über alle Tasks hinweg gilt: Verified, Implemented\textsubscript{static} und Attempted sind empirisch unterscheidbare Statuskategorien und fallen nicht notwendigerweise zusammen. Für die Trichterdarstellung wird Implemented$^\ast$ (Attempted $\wedge$ (Implemented\textsubscript{static} $\vee$ Verified)) als konsistente Evidenzstufe genutzt.
Kapitel~6 beschreibt diese Statusausweise als Ergebnisvariablen der Pipeline und nutzt sie als gemeinsame Referenz, um Code-, Test- und Chat-Artefakte pro Branch in Beziehung zu setzen.

\paragraph{Technische Qualität (aggregiert, ohne Kausalannahmen).}
Die in Kapitel~6 berichteten Qualitäts- und Technikindikatoren zeigen insgesamt eine heterogene Ausprägung über die Branches hinweg.
Dazu gehören (a) Unterschiede im Linting-Status und in der Häufigkeit typischer TypeScript-Surrogate (z.\,B. Any-Nutzung), (b) Spannweiten bei strukturbezogenen Metriken (z.\,B. Anzahl TypeScript-Dateien) sowie (c) variierende Testausgänge bis hin zu Konstellationen mit Fehlanteilen.
Diese Befunde werden in Kapitel~6 als Ko-Existenz von Pipeline-Fortschritt und technischen Merkmalen beschrieben. Es wird keine Korrelation oder Wirkbeziehung behauptet.
Zusätzlich wird sichtbar, dass sich technische Ausprägungen nicht auf einen einheitlichen "Projektzustand" reduzieren lassen: Branches mit ähnlichem Pipeline-Fortschritt können unterschiedliche Muster in Linting-, Typisierungs- und Testindikatoren aufweisen.
Zugleich werden die Messgrenzen betont: Die verwendeten Indikatoren ersetzen keine Coverage-, Komplexitäts- oder Performance-Metriken und erlauben entsprechend keine Aussagen zu nicht gemessenen Qualitätsdimensionen.

\paragraph{Chat-Interaktionsmuster (kompakt).}
Die Chat-Analysen zeigen große Spannweiten bei der Interaktionsintensität, insbesondere bei Wortanzahl und Prompt-Antwort-Zyklen.
Über die Branches hinweg koexistieren kurze, wenig strukturierte Prompts mit detaillierten Mehrpunkt-Prompts. Beide Formen treten im selben Datensatz nebeneinander auf.
Als wiederkehrendes Strukturmerkmal sind Debugging-Iterationen sichtbar (z.\,B. wiederholte Fehlermarker und erneute Anpassungsversuche), ohne dass daraus eine qualitative Einordnung ("gut/schlecht") abgeleitet wird.
Neben Intensitätsmaßen werden in Kapitel~6 auch Muster in der Interaktionsstruktur beschrieben, etwa wiederkehrende Task-Referenzen, Wechsel zwischen Aufgaben sowie Tooling-/Aktionsformulierungen im Verlauf der Chats.

\paragraph{Triangulation auf Meta-Ebene.}
Die Triangulation macht sichtbar, dass unterschiedliche Datenquellen teils kongruente, teils divergente Signale liefern: Pipeline-Status, technische Indikatoren, Testausgänge, Chat-Merkmale und Survey-Merkmale können für dieselben Branches zusammenpassen oder auseinanderlaufen.
Insbesondere verdeutlicht die Gegenüberstellung, dass Verified $\neq$ Implemented\textsubscript{static} $\neq$ Attempted und dass Spannungsfelder zwischen Statusausweisen, technischen Beobachtungen und Interaktionsmustern empirisch auftreten.
Auch Survey-Merkmale werden als kategoriale Kontextvariablen einbezogen (Selbsteinstufungsgruppen und benannte Hürdenkategorien), ohne dass daraus individuelle Zuschreibungen oder Bewertungen abgeleitet werden.
In der Gesamtschau entstehen damit nebeneinander stehende Ergebnislinien, die Gemeinsamkeiten und Unterschiede zwischen Branches sichtbar machen, ohne diese zu erklären.
Kapitel~6 stellt diese Spannungsfelder dar, löst sie jedoch nicht auf.

\paragraph{Abgrenzung zu Kapitel~7.}
Kapitel~6 beschreibt, \emph{was} beobachtet wurde. Kapitel~7 diskutiert, \emph{wie} diese Befunde einzuordnen sind und \emph{warum} bestimmte Muster auftreten könnten.


\chapter{Diskussion}
\label{chap:discussion}

\section{Ziel und Abgrenzung der Diskussion}
\label{sec:discussion_scope}

Kapitel~6 hat die Ergebnisse der Studie ausschließlich deskriptiv berichtet und die beobachteten Muster entlang mehrerer Datenquellen nebeneinandergestellt.
Kapitel~7 interpretiert diese Befunde im Licht des in Kapitel~2 eingeführten theoretischen Rahmens.
Ziel ist es, die Forschungsfragen argumentativ zu beantworten und die Hypothesen zu bewerten.
Damit adressiert die Diskussion die in Abschnitt~\ref{sec:related-work} formulierte Forschungslücke, indem Befunde aus einem kontrollierten Brownfield-Setting entlang der Pipeline-Status und im Kompetenzkontext eingeordnet werden, ohne daraus eine Generalisierbarkeit über andere Settings oder KI-Assistenzsysteme abzuleiten.
Dabei werden keine neuen Daten eingeführt.
Es werden keine zusätzlichen Metriken berechnet.
Es werden keine Kausalzusammenhänge behauptet, die durch das Studiendesign oder die erhobenen Artefakte nicht abgesichert sind.

\paragraph{Methodische Selbstbegrenzung (Kausalität, Gruppenvergleiche, Kompetenzgruppen).}
Die Studie ist nicht als kausales Design angelegt, sondern als kontrollierte, artefaktbasierte Beobachtungsstudie: Es gibt keine Randomisierung, keine experimentelle Manipulation und keine prozessnahen Messungen, die eine Identifikation von Ursache--Wirkung erlauben würden.
Zudem sind zentrale Einflussgrößen (z.B. Technologie-Fit, Test-/Debugging-Routinen, Prompting-Erfahrung, Unterbrechungen) nur teilweise messbar und können daher nicht hinreichend kontrolliert werden.
Aus diesem Grund werden Befunde konsistent als deskriptive Muster und theoriegeleitete Einordnungen berichtet.
Auf inferenzstatistische Gruppenvergleiche (Signifikanztests, Effektstärken) wird verzichtet, da die Fallzahlen pro Kompetenzgruppe klein sind, die Messannahmen für stabile Schätzungen nicht belastbar geprüft werden können und multiple Konfundierungen das Risiko scheinbarer Gruppenunterschiede erhöhen.
Die Kompetenzgruppen werden deshalb nicht als Effektvariablen im Sinne gruppenstabiler Leistungs- oder Qualitätsunterschiede modelliert, sondern als Analyseperspektiven zur Strukturierung von Nutzungstypen und Interpretationskontexten.
Damit bleibt die Diskussion auf die Frage fokussiert, welche Artefaktspuren und Prozesssignale im Brownfield-Setting beobachtbar sind und wie diese im Kompetenzkontext plausibel eingeordnet werden können.

\paragraph{Geltungsbereich der Ergebnisse.}
Die in dieser Arbeit diskutierten Befunde beziehen sich ausschließlich auf \emph{interaktive} KI-Assistenzsysteme im Vibe-Coding-Kontext, die in einer Brownfield-Sandbox auf ein bestehendes Codeartefakt angewendet werden. Sie gelten für eine zeitlich begrenzte Aufgabenbearbeitung unter den in Kapitel~6 beschriebenen Rahmenbedingungen. Nicht Gegenstand der Aussagen sind Greenfield-Projekte oder langfristige Produktentwicklung. Ebenfalls nicht abgedeckt sind autonome Agenten mit eigener Aufgabenplanung und eigenständiger Ziel- und Teilzielbildung.

Eine zentrale Linse der Diskussion sind die vier Kompetenzgruppen No-Code, Low-Code, Mid-Code und High-Code.
Diese Gruppen werden als gleichwertige Analyseperspektiven behandelt.
Sie sind keine Qualitätsstufen und dienen nicht der Bewertung einzelner Personen.
Aussagen beziehen sich daher auf aggregierte Gruppenmuster oder auf einzelne Branches als Analyseeinheiten.
Eine Leistungsbewertung oder Rangordnung einzelner Personen erfolgt nicht.

\paragraph{Group comparability \& confounders.}
Die Vergleichbarkeit der Kompetenzgruppen ist in dieser Studie nur in Teilen abgesichert. Auf der einen Seite sind zentrale Rahmenbedingungen für alle Teilnehmenden konstant gehalten (einheitliche IDE \enquote{VS Code}, identische Sandbox und Aufgabenstellung, identisches Zeitlimit, sowie dasselbe KI-Assistenzsystem aus GitHub Copilot und Claude Sonnet~4.5). Damit sind Unterschiede in den in Kapitel~6 berichteten Artefaktspuren und Pipeline-Status grundsätzlich \emph{kompatibel mit} kompetenzabhängigen Nutzungstypen. Auf der anderen Seite können alternative Erklärungen nicht ausgeschlossen werden: Unterschiede in Technologie-Fit (TS/Node/Backend), allgemeiner Berufserfahrung, Test-/Debugging-Routinen und individueller Risikotoleranz beeinflussen sowohl Implementationsentscheidungen als auch die Wahrscheinlichkeit, \textit{Verified} unter der gegebenen Testmethodik zu erreichen. Einige dieser Faktoren sind im Pre-Survey als Kontextmerkmale verfügbar (z.B. Erfahrung, Technologie-Fit) und können deskriptiv als Gruppenprofil berichtet werden; andere (z.B. Vorerfahrung mit Copilot/LLMs, Prompting-Strategiequalität, Hardware/Performance, Unterbrechungen) sind in den vorliegenden Artefakten nicht belastbar operationalisiert. Entsprechend sind die Diskussionen in Kapitel~7 als theoriegeleitete Einordnungen zu lesen: Unterschiede \emph{lassen sich interpretieren als} unterschiedliche Integrations- und Validierungsmodi im Kompetenzkontext, ohne daraus kausale Effekte oder gruppenstabile Leistungsrangfolgen abzuleiten.

Zur begrifflichen Trennung wird -- wie in Kapitel~1 eingeführt -- zwischen Assistenzumgebung (Tool), Modell (LLM) und \emph{KI-Assistenzsystem} als funktionaler Einheit unterschieden.

\section{Rückbindung an die Hypothesen}
\label{sec:discussion_hypotheses}

Die Hypothesen wurden im Forschungsdesign als Erwartung darüber formuliert, wie Kompetenz und Arbeitsmodus mit dem KI-Assistenzsystem zusammenwirken.
Im Folgenden wird jede Hypothese anhand der in Kapitel~6 berichteten Befundlinien diskutiert.
Dabei werden die vier Kompetenzgruppen No-Code, Low-Code, Mid-Code und High-Code als Analyseperspektiven verwendet.
Nicht jede Hypothese zielt in gleicher Weise auf alle Gruppen.

\subsection{H1: Mid-Code als Sweet Spot}
\label{sec:discussion_h1}

\textbf{Befunde, die H1 stützen.}
Die Ergebnisse aus Kapitel~6 sprechen für ein Muster, das mit einer produktiven Balance zwischen Delegation und Kontrolle bei Mid-Code vereinbar ist.
In der Triangulation erscheinen für Mid-Code sowohl substanzielle Umsetzungsstände als auch Hinweise auf systematische Validierung als plausibles Arbeitsmuster.
Dieses Bild ist kompatibel mit der in Kapitel~2 beschriebenen Kompetenzlogik, nach der Mid-Code Probleme in Teilprobleme zerlegt und die durch das KI-Assistenzsystem erzeugten prozeduralen Chunks fachlich prüfen kann.
Die in Kapitel~6 beobachtete Koexistenz von Implementationsartefakten und Debugging-Iterationen ist damit konsistent.

\textbf{Befunde, die H1 relativieren.}
Die Gegenüberstellung Kompetenzgruppe und Umsetzungsstand zeigt kein exklusives Muster, das nur Mid-Code zugeordnet werden kann.
Auch in anderen Gruppen treten Branches mit weiter reichenden Umsetzungsständen auf.
Damit ist das empirische Signal nicht stark genug, um Mid-Code als eindeutig vorteilhaften Nutzungstyp gegenüber den anderen Analyseperspektiven zu interpretieren.
Zudem sind die in Kapitel~6 ausgewiesenen Qualitätsindikatoren nur Surrogate.
Sie erlauben keine direkte Aussage darüber, ob Mid-Code tatsächlich systematischer oder effizienter validiert.

\textbf{Bewertung.}
H1 wird insgesamt teilweise gestützt.
Die Befunde sind kompatibel mit einem Sweet-Spot-Argument.
Sie sind jedoch nicht hinreichend spezifisch, um eine klare Stützung gegenüber den anderen Kompetenzgruppen zu begründen.
Damit kann die Arbeit H1 als theoriekompatibles Muster im Brownfield-Setting auf Ebene von Pipeline- und Artefaktspuren einordnen; offen bleibt, ob und unter welchen Variationen von Aufgaben, Zeitbudget und prozessnahen Messungen ein solcher Sweet Spot replizierbar und trennscharf nachweisbar ist.

\subsection{H2: Validierungsparadoxon bei Low-Code}
\label{sec:discussion_h2}

\textbf{Befunde, die H2 stützen.}
Kapitel~6 zeigt wiederholt Spannungsfelder zwischen sichtbarer Implementationsaktivität und fehlender testbasierter Bestätigung, insbesondere bei konfigurationssensitiven und stark integrierten Aufgaben.
Die Befunde sind vereinbar mit der Interpretation, dass eine rasche Generierung von Artefakten nicht automatisch mit stabiler, testkonformer Funktionalität einhergeht.
Dies ist kompatibel mit Kapitel~2, in dem Vibe-Coding als System~1-Aktivität und Automation Bias als Risiko unkritischer Übernahme diskutiert wird.
Für Low-Code lässt sich dies theoriegeleitet damit plausibilisieren, dass syntaktische Plausibilität häufig leichter zu beurteilen ist als semantische Korrektheit in einem bestehenden System.

\textbf{Befunde, die H2 relativieren.}
Die Artefakte in Kapitel~6 erlauben keine direkte Beobachtung individueller Validierungspraktiken, etwa durch Protokolle von Testläufen oder IDE-Aktionen.
Auch werden Chat- und Code-Metriken nicht systematisch nach Kompetenzgruppe aufgeschlüsselt.
Damit bleibt die Zuordnung des Validierungsparadoxons zu Low-Code eine theoriegeleitete Interpretation.
Es handelt sich nicht um einen unmittelbar gruppenspezifisch gemessenen Effekt.

\textbf{Bewertung.}
H2 ist plausibel und wird durch die in Kapitel~6 berichteten Diskrepanzen zwischen Statuskategorien inhaltlich gestützt.
Aufgrund der fehlenden gruppenspezifischen Validierungsdaten bleibt die Hypothese jedoch nicht eindeutig.
Damit kann die Arbeit die analytische Relevanz der Trennung von Implementationsartefakten und testgebundener Bestätigung für H2 sichtbar machen; offen bleibt, ob und in welchem Ausmaß das Muster in prozessnahen Daten tatsächlich kompetenzspezifisch durch konkrete Validierungspraktiken getrieben ist.

\subsection{H3: Skepsis und Qualitätsfokus bei High-Code}
\label{sec:discussion_h3}

\textbf{Befunde, die H3 stützen.}
Kapitel~6 liefert Hinweise, die mit der Beobachtung vereinbar sind, dass anspruchsvolle Integrationsaufgaben unter Zeitdruck häufig nicht bis zur testbasierten Stabilisierung gelangen.
Für High-Code ist im Rahmen von Kapitel~2 ein Arbeitsmodus plausibel, der Vorschläge stärker hinterfragt und damit mehr kognitive Ressourcen in Kontrolle investiert.
Dieses Muster wäre vereinbar mit kalibriertem Vertrauen, bei dem Overtrust vermieden wird.
Es wäre ebenfalls vereinbar mit dem in Kapitel~2 diskutierten Befund, dass das Modifizieren von KI-Code kognitiv teuer sein kann.

\textbf{Befunde, die H3 relativieren.}
Die in Kapitel~6 berichteten Chat-Muster und Qualitätsindikatoren werden nicht nach Kompetenzgruppen getrennt.
Daraus folgt, dass ein erhöhter Qualitätsfokus von High-Code nicht direkt beobachtet, sondern nur aus der Theorie abgeleitet werden kann.
Zudem zeigen die Ergebnisse keine eindeutige Dominanz von High-Code beim Umsetzungsstand.
Ein höherer Qualitätsanspruch kann unter Zeitlimit sogar zu konservativen Entscheidungen führen, ohne dass dies als geringere Kompetenz zu interpretieren ist.

\textbf{Bewertung.}
H3 ist im theoretischen Rahmen gut begründbar.
Die empirischen Befunde aus Kapitel~6 sind mit dieser Hypothese vereinbar, stützen sie jedoch nicht eindeutig.
Damit kann die Arbeit H3 als vorsichtige, theoriegeleitete Erklärung für mögliche Kontroll- und Integrationsmodi im Zeitlimit diskutieren; offen bleibt, ob sich ein entsprechender Qualitätsfokus in prozessnahen Messungen (z.B. Review-Tiefe, Testlauf-Sequenzen) als distinktes Muster empirisch abgrenzen lässt.

\section{Kompetenzabhängige Interaktionsmuster}
\label{sec:discussion_competence_patterns}

Kapitel~6 zeigt, dass Branches sich nicht nur im Umsetzungsstand unterscheiden, sondern auch in der Art der Interaktion, in der Struktur der Prompts und in der Häufigkeit von Debugging-Sequenzen.
Diese Variation lässt sich im Licht von Kapitel~2 als Ausdruck unterschiedlicher Kompetenzlagen interpretieren.
Die folgenden Deutungen sind explizit theoriegeleitet und formulieren keinen kausalen Anspruch.
Sie beschreiben Nutzungstypen und Analyseperspektiven, nicht Qualitätsstufen oder individuelle Leistungsfähigkeit.
Zentral ist dabei nicht das Tool selbst, da alle Teilnehmenden mit derselben Assistenzumgebung (GitHub Copilot) und derselben Modellkonfiguration gearbeitet haben.
In dieser Einordnung ist vielmehr relevant, welche kognitiven Ressourcen und welche mentalen Modelle zur Verfügung stehen, um Vorschläge zu prüfen, zu integrieren und bei Bedarf zu korrigieren.

\subsection{No-Code}
\label{sec:discussion_no_code}

Die Befunde aus Kapitel~6 lassen sich so lesen, dass für No-Code der Zugang zu umsetzungsnahen Artefakten erleichtert werden kann.
Gleichzeitig kann Validierung und Integration im bestehenden System häufig anspruchsvoll bleiben.
Im Rahmen von Kapitel~2 lässt sich dies als Vibe-Coding-Arbeitsmodus interpretieren, in dem das KI-Assistenzsystem prozedurale Chunks liefert, die fehlende Produktionsregeln überbrücken.
Gleichzeitig kann ohne tragfähige Schemata der Intrinsic Load der Prüfung und Integration höher ausfallen.
Damit ist es im Sinne von Kapitel~2 plausibel, dass System~2 situativ nicht in dem Maß aktiviert wird, das für semantische Validierung erforderlich wäre, obwohl die Oberfläche plausibel wirkt.
Diese Einordnung bleibt theoriegeleitet, da Kapitel~6 keine direkten Prozessmessungen der Validierung auf individueller Ebene enthält.

\subsection{Low-Code}
\label{sec:discussion_low_code}

Kapitel~6 lässt sich so lesen, dass Low-Code sichtbare Implementationsartefakte erzeugt, während testgebundene Bestätigung bei stärker integrierten Aufgaben nicht durchgängig erreicht wird.
Im Rahmen von Kapitel~2 lässt sich dies als Spannungsfeld zwischen Vertrauen und Kontrolle interpretieren.
Dabei kann das KI-Assistenzsystem als Beschleuniger genutzt werden, während semantische Korrektheit im Systemkontext schwerer zu prüfen bleibt.
Wenn Validierungsroutinen nicht etabliert sind, ist in dieser theoretischen Lesart plausibel, dass unkalibriertes Vertrauen häufiger auftritt, ohne dass dies aus den Artefakten als gruppenspezifischer Effekt gemessen werden kann.
Die in Kapitel~6 berichteten Strategiemuster, etwa tutorial-nahe Implementationsflüsse und Boilerplate, sind mit einem stärker template-basierten Arbeitsmodus vereinbar.
Diese Einordnung bleibt theoriegeleitet, da Kapitel~6 Validierungshandlungen nicht gruppenspezifisch und prozessnah erfasst.

\subsection{Mid-Code}
\label{sec:discussion_mid_code}

Kapitel~6 ist mit der Beobachtung kompatibel, dass Mid-Code sowohl Umsetzungsartefakte als auch Iterationen der Nachjustierung und Validierung zeigt.
Im Rahmen von Kapitel~2 lässt sich dies als Kompetenzlage interpretieren, die zielorientierte Zerlegung und aktive Kontrolle kombiniert.
Damit lässt sich das KI-Assistenzsystem eher als Werkzeug in einem kontrollierten Problemlöseprozess einordnen.
Die in Kapitel~6 berichtete Vielfalt an Prompt-Formen ist als funktionale Variation plausibel, in der kurze Prompts Beschleunigung und strukturierte Prompts Spezifikation und Kontextabgleich unterstützen.
Debugging-Iterationen sind dabei als normaler Bestandteil iterativer Integrationsarbeit zu lesen, nicht als Effizienzindikator.
Diese Einordnung bleibt vorsichtig, da Kapitel~6 keine direkten Messungen von Review- und Validierungstiefe enthält.

\subsection{High-Code}
\label{sec:discussion_high_code}

Kapitel~6 ist kompatibel mit dem Befund, dass unter Zeitlimit bei anspruchsvollen Integrationsaufgaben testgebundene Bestätigung nicht durchgängig erreicht wird.
Im Rahmen von Kapitel~2 lässt sich dies für High-Code als Nutzungstyp interpretieren.
In dieser Einordnung wird das KI-Assistenzsystem eher für Boilerplate und Exploration genutzt, während zentrale Entscheidungen und Integrationslogik stärker manuell geprüft werden.
In dieser Perspektive können intensives Prüfen und konservative Integrationsentscheidungen mit geringerem sichtbarem Fortschritt koexistieren, ohne dass daraus unmittelbar Rückschlüsse auf Qualität oder Kompetenz abzuleiten wären.
Diese Einordnung bleibt vorsichtig, da Kapitel~6 keine direkte Messung von Review-Tiefe oder Kontrollhandlungen enthält.

\section{Interpretation der Pipeline-Muster}
\label{sec:discussion_pipeline}

Die Pipeline-Kategorien \textit{Attempted}, \textit{Implemented\textsubscript{static}} und \textit{Verified} strukturieren die Befunde in Kapitel~6 entlang eines Trichters; für die Trichterdarstellung wird dort zusätzlich \textit{Implemented$^\ast$} als abgeleitete Evidenzstufe verwendet.
Die folgenden Deutungen binden diese deskriptiven Status-Diskrepanzen explizit an die in Kapitel~2 eingeführten Konzepte (u.a. Cognitive Load, Automation Bias) zurück.
Für die Diskussion ist entscheidend, dass diese Kategorien unterschiedliche Aspekte desselben Arbeitsprozesses repräsentieren.
In der Notation gilt: $\textit{Attempted} \neq \textit{Implemented\textsubscript{static}} \neq \textit{Verified}$.
Attempted beschreibt Intention und Planbildung auf der Chat-Ebene.
Implemented\textsubscript{static} beschreibt aufgabenspezifische Artefakte in der Codebasis (White-Box).
Verified beschreibt testgebundenes, von außen beobachtbares Verhalten (Black-Box).
Verified ist damit methodengebunden und als Integrationsindikator zu lesen, nicht als Leistungs- oder Qualitätsmaß.
Aus Kapitel~6 folgt, dass diese Ebenen nicht zusammenfallen müssen.

\subsection{Attempted ist nicht gleich \texorpdfstring{Implemented\textsubscript{static}}{Implementedstatic}}
\label{sec:discussion_attempted_vs_implemented}

Die beobachtete Differenz zwischen Attempted und Implemented\textsubscript{static} lässt sich plausibel als Integrationsphänomen interpretieren.
Insbesondere Aufgaben, die mehrere Systemteile koppeln, gehen mit zusätzlichem Koordinationsaufwand einher.
Im Rahmen von Kapitel~2 lässt sich dies über Cognitive Load erklären.
Die Aufgabe reduziert sich nicht auf das Erzeugen von Code, sondern umfasst das Verstehen des bestehenden Kontexts, das Abstimmen von Datenmodellen und Schnittstellen sowie das Beheben von Fehlanpassungen.
Dieser Aufwand dürfte den Intrinsic Load erhöhen.
Das KI-Assistenzsystem kann Extraneous Load senken, indem es Syntax und Boilerplate generiert.
Es kann die semantische Konsistenzarbeit im Systemkontext jedoch nicht ersetzen.

\subsection{\texorpdfstring{Implemented\textsubscript{static}}{Implementedstatic} ist nicht gleich Verified}
\label{sec:discussion_implemented_vs_verified}

Die Diskrepanz zwischen Implemented\textsubscript{static} und Verified ist nicht als Qualitätsurteil über einzelne Branches zu lesen.
Kapitel~6 betont, dass Verified ein Diagnosesignal innerhalb der eingesetzten Testmethodik ist.
Für konfigurationssensitive Aufgaben kann Verified stark von Details des Endpunktvertrags und der Testannahmen abhängen.
Ein fehlendes Verified-Signal kann daher bedeuten, dass eine Lösung heterogen, nur teilweise integriert oder nicht exakt testkonform ist.
Es kann auch bedeuten, dass die Tests eine bestimmte Implementationsform voraussetzen, die nicht mit allen plausiblen Lösungen kompatibel ist.

\subsection{Warum Verified gleich null kein eindeutiges Negativsignal ist}
\label{sec:discussion_verified_zero}

Kapitel~6 zeigt für bestimmte Integrationsaufgaben, dass testbasierte Bestätigung ausbleibt.
Im Kontext dieser Studie ist das als erwartbares Ergebnis eines strikten Zeitlimits und hoher Schnittstellenkomplexität interpretierbar.
Gerade in einer Brownfield-Umgebung dürfte der Aufwand durch bestehende Konventionen, vorhandene Datenmodelle und implizite Vertragsannahmen steigen.
Dies kann mit erhöhter kognitiver Belastung und Debugging-Aufwand einhergehen.
Unter Automation Bias Perspektive ist zudem plausibel, dass scheinbar fertige Artefakte in einzelnen Fällen zu früh als ausreichend akzeptiert werden.
Die Abwesenheit von Verified ist in dieser Arbeit daher kein Synonym für fehlendes Verständnis, sondern kann als Signal für Integrationskosten und Validierungsengpässe gelesen werden.

\section{Einordnung der Chat-Interaktionsmuster}
\label{sec:discussion_chat}

Die Chat-Ergebnisse aus Kapitel~6 zeigen starke Variation der Interaktionsintensität und wiederkehrende Muster von Debugging-Sequenzen.
Kapitel~6 bewertet diese Muster nicht.
Für die Diskussion sind drei Punkte zentral.

\subsection{Warum Chat-Intensität kein Effizienzmaß ist}
\label{sec:discussion_chat_intensity}

Eine hohe Interaktionsintensität kann mehrere Bedeutungen haben.
Sie kann auf Exploration, auf unklare Anforderungen, auf Fehlanpassungen an den Kontext oder auf konsequente Validierung und Nachjustierung hinweisen.
Kapitel~2 betont, dass Vibe-Coding den Schwerpunkt vom Schreiben zum Lesen und Prüfen verschiebt.
Damit kann ein längerer Chat-Verlauf auch Ausdruck von Kontrolle und Review sein.
Umgekehrt kann ein kurzer Chat-Verlauf sowohl effiziente Beschleunigung als auch unkritische Übernahme bedeuten.
Ohne zusätzliche Prozessdaten, etwa Messungen von Testläufen oder IDE-Interaktionen, ist eine Effizienzdeutung nicht abgesichert.

\subsection{Warum Debugging-Loops normal sind}
\label{sec:discussion_debugging_loops}

Kapitel~6 berichtet Debugging-Iterationen als wiederkehrendes Strukturmerkmal.
Im Rahmen von Kapitel~2 ist dies erwartbar.
Mixed-Initiative Interaction bedeutet, dass Vorschläge, Rückfragen und Korrekturen in Schleifen organisiert sind.
Zudem kann die Arbeit in einem bestehenden System die Wahrscheinlichkeit von Konfigurations- und Schnittstellenfehlern erhöhen.
Debugging-Sequenzen sind daher als normaler Bestandteil von Integrationsarbeit zu interpretieren.
Sie sind nicht per se ein Indikator für geringe Kompetenz oder geringe Qualität.

\subsection{Explorative und zielgerichtete Interaktion}
\label{sec:discussion_exploration_vs_acceleration}

Kapitel~2 unterscheidet Interaktionsmodi wie Exploration und Beschleunigung.
Kapitel~6 zeigt sowohl kurze, wenig strukturierte Prompts als auch strukturierte, mehrpunktige Prompts.
Diese Koexistenz lässt sich als situativ wechselnder Arbeitsmodus interpretieren.
Exploration kann genutzt werden, um APIs, Projektkonventionen oder mögliche Lösungswege zu verstehen.
Beschleunigung kann genutzt werden, um bekannte Muster schnell zu implementieren.
Der Wechsel zwischen diesen Modi ist in einer zeitlich begrenzten Aufgabenbearbeitung plausibel.

\section{Implikationen für Praxis und Forschung}
\label{sec:discussion_implications}

Aus der Diskussion lassen sich vorsichtige Implikationen für Praxis und Forschung formulieren, die über die konkrete Sandbox hinausreichen.
Sie betreffen sowohl die Einführung von KI-Tools in Organisationen als auch die Art, wie deren Nutzen evaluiert wird.
Sie knüpfen an die in Kapitel~6 berichteten Befunde an und werden im Licht der in Kapitel~2 eingeführten Konzepte formuliert.

\subsection{Implikationen für Unternehmen}
\label{sec:discussion_implications_practice}

Die Bereitstellung eines KI-Tools allein adressiert nicht automatisch die in Kapitel~6 sichtbaren Unterschiede in Nutzungstypen.
Die Befunde sind kompatibel mit der Annahme, dass Kompetenz den Umgang mit Delegation, Kontrolle und Validierung mitprägt.
Als Implikation lässt sich vorsichtig ableiten, dass organisatorische Einbettung und Qualifizierung den beobachteten Prozessunterschieden begegnen können.
Dies kann sich etwa in Kompetenzaufbau zu Validierung, Testdenken und Schnittstellenverträgen sowie in expliziten Erwartungen an Review und Integrationsarbeit ausdrücken.

\subsection{Implikationen für die Evaluation von KI-Coding-Tools}
\label{sec:discussion_implications_evaluation}

Kapitel~6 zeigt, dass einzelne Signale wie Teststatus, Linting oder Interaktionsintensität isoliert schwer interpretierbar sind.
Für Forschung und Praxis ergibt sich daraus, dass Interpretationen belastbarer werden, wenn mehrere Perspektiven gemeinsam betrachtet werden.
Insbesondere ist die Unterscheidung zwischen Intention, Implementationsartefakt und testbasierter Bestätigung zentral.
Ebenso ist der Kompetenzkontext für die Einordnung der Artefakte relevant, weil aggregierte Durchschnittseffekte Nutzungstypen überdecken können.

\subsection{Abgrenzung zu Marketing-Narrativen}
\label{sec:discussion_marketing}

In der öffentlichen Darstellung werden häufig generische Produktivitätsgewinne hervorgehoben.
Die in Kapitel~6 berichteten Muster legen dagegen eine differenzierte Einordnung entlang von Integrationsaufwand, Validierbarkeit und Kompetenzkontext nahe.
Im Rahmen von Kapitel~2 können KI-Assistenzsysteme Extraneous Load reduzieren und damit Umsetzungsschritte beschleunigen.
Gleichzeitig verschieben sich Anforderungen hin zu Review, semantischer Validierung und testmethodisch gebundener Absicherung.
Ohne diese Prozessanteile ist im Sinne von Kapitel~2 unkalibriertes Vertrauen plausibel, ohne dass dies kausal belegt wäre.
Damit ist der Nutzen nicht als automatische Eigenschaft eines Tools zu verstehen, sondern als Ergebnis einer arbeitsorganisatorischen Einbettung, die Kompetenzlagen und Integrationskosten berücksichtigt.

\section{Limitationen und Ausblick}
\label{sec:discussion_limitations}

Die Studie ist in ihrer Aussagekraft durch mehrere Faktoren begrenzt.
Erstens ist die Stichprobengröße begrenzt, wodurch gruppenspezifische Effekte nur eingeschränkt trennscharf beobachtbar sind.
Zweitens setzt das strikte Zeitlimit den Fokus auf frühe Integrationsschritte und kann unvollständige Stabilisierung begünstigen.
Drittens wurde GitHub Copilot als Assistenzumgebung in Kombination mit dem Modell Claude Sonnet~4.5 betrachtet, wodurch sich keine Aussagen zur Generalisierbarkeit über andere KI-Assistenzsysteme oder Modelle ableiten lassen.
Viertens liegt keine Langzeitbeobachtung vor, sodass Lern- und Adaptationseffekte über Wochen oder Monate nicht erfasst werden.

Aus diesen Limitationen ergeben sich unmittelbare Ansatzpunkte für weitere Forschung.
Sinnvoll sind längere Beobachtungsfenster mit wiederholten Sessions, um Stabilisierung, Refactoring und Validierungsroutinen erfassen zu können.
Ebenso sind Designs relevant, die Prozessdaten stärker erfassen, etwa Testlauf-Sequenzen, Review-Handlungen oder Kontextbereitstellung im Editor.
Schließlich ist es für künftige Arbeiten relevant zu prüfen, inwieweit die beobachteten Muster über unterschiedliche KI-Assistenzsysteme und Modellkonfigurationen hinweg stabil sind, ohne daraus unmittelbare Leistungsrangfolgen abzuleiten.

Kapitel~8 fasst die zentralen Erkenntnisse zusammen und leitet daraus die abschließenden Antworten auf die Forschungsfragen ab.

\chapter{Fazit und Schlussfolgerungen}
\label{chap:fazit}

\section{Ziel und Einordnung des Kapitels}
\label{sec:conclusion_scope}

Kapitel~8 schließt die Arbeit ab, indem die in Kapitel~6 berichteten Ergebnisse zusammengeführt und die in Kapitel~7 diskutierten Einordnungen in prägnante Schlussfolgerungen überführt werden.
Im Unterschied zu Kapitel~6 (deskriptiv) und Kapitel~7 (theoriegeleitet) dient dieses Kapitel der abschließenden Synthese und der expliziten Beantwortung der Forschungsfragen.
Es werden keine neuen Daten eingeführt und keine zusätzlichen Metriken berechnet.

Übergeordnetes Ziel der Arbeit war die Evaluation eines interaktiven KI-Assistenzsystems im Vibe-Coding-Kontext in einer kontrollierten Brownfield-Sandbox.
Im Studiensetup wurde -- konsistent zur Begriffsklärung in Kapitel~1 -- zwischen Assistenzumgebung (Tool), Modell (LLM) und KI-Assistenzsystem als funktionaler Einheit unterschieden.

\section{Beantwortung der Forschungsfragen}
\label{sec:conclusion_rqs}

\subsection[Forschungsfrage 1: Kompetenz und Pipeline-Status]{Forschungsfrage 1: Wie hängt die Entwicklerkompetenz mit dem Erreichen der Pipeline-Status Attempted, \texorpdfstring{Implemented\textsubscript{static}}{Implementedstatic} und Verified zusammen?}
\label{sec:conclusion_rq1}

\textbf{Kurzantwort.}
Attempted und Implemented\textsubscript{static} sind über die Kompetenzgruppen hinweg als Artefaktspuren sichtbar; Verified als testgebundene Stufe wird am seltensten erreicht.

\textbf{Verdichtete Begründung.}
Kapitel~6 trennt die Evidenzebenen Attempted, Implemented\textsubscript{static} und Verified.
Attempted steht für Intention (Chat-Ebene), Implemented\textsubscript{static} für White-Box-Artefakte (Codebasis), Verified für testgebundenes Black-Box-Verhalten und ist damit methodengebunden.
Für die Trichterdarstellung wird in Kapitel~6 zusätzlich Implemented$^\ast$ als abgeleitete Evidenzstufe verwendet.
Die Befunde aus Kapitel~6 und die Einordnung in Kapitel~7 sind damit vereinbar, dass Implementationsfortschritt nicht notwendigerweise in testgebundene Bestätigung übergeht.
Kompetenz dient dabei als Analyseperspektive auf unterschiedliche Integrations- und Validierungsmodi, nicht als Leistungsbewertung.

\subsection[Forschungsfrage 2: Code- und Test-Artefakte]{Forschungsfrage 2: Welche Rolle spielen Code- und Test-Artefakte für die Bewertung der Arbeitsergebnisse eines KI-Assistenzsystems?}
\label{sec:conclusion_rq2}

\textbf{Kurzantwort.}
Code- und Test-Artefakte sind zentral, weil sie die Trennung zwischen sichtbarer Implementierung (Implemented\textsubscript{static}) und testgebundener Bestätigung (Verified) operationalisieren und Chat-Signale methodisch ergänzen.

\textbf{Verdichtete Begründung.}
Die Artefaktanalyse macht Implementationsfortschritt über Quellcodeänderungen, Surrogat-Signale (z.B. Linting) und strukturelle Spuren in der Codebasis nachvollziehbar.
Testgebundene Ergebnisse bilden beobachtbares Verhalten unter den gegebenen Tests ab und sind damit zentral für die Bestätigung im Systemkontext.
Die Pipeline-Logik liefert so eine nachvollziehbare Trennung von Intention, Implementationsartefakt und testgebundener Bestätigung (Kapitel~6) und wird in Kapitel~7 methodisch eingeordnet.
Kompetenz ergänzt diese Sicht als Analyseperspektive, weil sich Artefaktspuren nicht sinnvoll in einem einzelnen, gruppenunabhängigen Signal verdichten lassen.

\subsection[Forschungsfrage 3: Chat-Interaktionsmuster]{Forschungsfrage 3: Welche Chat-Interaktionsmuster treten im Vibe-Coding-Kontext auf, und wie sind sie für die Evaluation einzuordnen?}
\label{sec:conclusion_rq3}

\textbf{Kurzantwort.}
Im Vibe-Coding-Kontext variieren Chat-Interaktionsmuster deutlich; ihre Intensität ist weder als Effizienzmaß noch als Qualitätsmaß geeignet.
Für die Evaluation sind Chat-Muster nur im Zusammenspiel mit Pipeline-Status sowie Code- und Test-Artefakten aussagekräftig.

\textbf{Verdichtete Begründung.}
Die Chat-Artefakte zeigen unterschiedliche Prompt-Strukturen, wechselnde Interaktionsmodi und wiederkehrende Debugging-Schleifen.
Diese Muster sind konsistent damit, dass Interaktion sowohl Exploration als auch Beschleunigung unterstützen kann.
Eine hohe Interaktionsintensität kann dabei gleichermaßen aus Validierung, aus Fehlanpassungen an den Systemkontext oder aus iterativer Integration resultieren.
Entsprechend ist die Einordnung von Chat-Mustern an die Pipeline-Logik zu koppeln.
Attempted macht die Intention sichtbar, Implemented\textsubscript{static} die resultierenden Artefakte und Verified die testgebundene Funktion im System.
Die Kompetenzgruppen liefern hierfür den Kontext als Analyseperspektiven auf Nutzungsformen, ohne dass daraus eine Bewertung einzelner Personen abgeleitet wird.

\section{Zentrale Erkenntnisse der Arbeit}
\label{sec:conclusion_key_findings}

Die Arbeit verdichtet sich in folgenden Befundlinien.

\begin{itemize}
    \item \textbf{Kompetenz als Perspektive auf Nutzungstypen:} Unterschiede zeigen sich in Nutzungsformen des KI-Assistenzsystems, nicht als Wertung von Personen oder Gruppen.
    \item \textbf{Trennung von Intention, Artefakt und Bestätigung:} Attempted, Implemented\textsubscript{static} und Verified bilden unterschiedliche Ebenen desselben Arbeitsprozesses ab und fallen nicht notwendigerweise zusammen.
    \item \textbf{Validierung und Integration als Engpass:} Sichtbare Implementationsartefakte sind häufig erreichbar, während testgebundene Bestätigung die methodisch engste Stufe darstellt.
    \item \textbf{Chat-Interaktion ist kein Ersatz für Ergebnisprüfung:} Interaktionsintensität und Prompt-Formen sind ohne Artefakt- und Testbezug nicht belastbar als Effizienz- oder Qualitätssignale.
    \item \textbf{Grenzen vereinfachter Produktivitätsnarrative:} Produktivität lässt sich in der untersuchten Umgebung nicht auf Generierungsgeschwindigkeit reduzieren, sondern ist entlang von Integrationsaufwand und Validierungsarbeit einzuordnen.
    \item \textbf{Relevanz von Triangulation statt Einzelmetriken:} Belastbare Einordnung entsteht durch das Zusammenspiel von Pipeline-Status, Code- und Test-Artefakten sowie Chat-Mustern.
\end{itemize}

\section{Beitrag der Arbeit}
\label{sec:conclusion_contributions}

\textbf{Wissenschaftlicher Beitrag.}
Die Arbeit liefert eine strukturierte Einordnung der Evaluation eines KI-Assistenzsystems im Vibe-Coding-Kontext, indem Kompetenzgruppen als Analyseperspektiven mit Artefaktspuren und Prozesssignalen zusammengeführt werden.

\textbf{Methodischer Beitrag.}
Die Arbeit operationalisiert Aufgabenerfolg über die Pipeline-Logik (Attempted, Implemented\textsubscript{static}, Verified; sowie Implemented$^\ast$ als abgeleitete Evidenzstufe) und verknüpft diese systematisch mit der Triangulation mehrerer Artefaktquellen.
Damit liegt ein Evaluationsschema vor, das die Trennung zwischen intendierter Bearbeitung, sichtbarer Implementierung und testgebundener Bestätigung explizit und nachvollziehbar macht.

\textbf{Konzeptioneller Beitrag.}
Die Arbeit betont Kompetenz als Perspektive auf Nutzungstypen und verschiebt den Fokus von einer rein toolzentrierten Betrachtung hin zu der Frage, wie KI-Assistenzsysteme in konkreten Entwicklungspraktiken genutzt und geprüft werden.

\textbf{Empirischer Beitrag.}
Die Sandbox-Studie in einer Brownfield-Umgebung liefert eine Artefaktbasis aus Code-, Test- und Chat-Spuren, anhand derer die Pipeline-Status operationalisiert und gruppenbezogen eingeordnet werden können (Kapitel~6/7).

\textbf{Praktischer Beitrag.}
Für Organisationen ergibt sich ein Rahmen zur Einführung und Evaluation von KI-Coding-Tools, der Qualifizierung, Review-Praxis und testbasierte Absicherung als Bestandteile einer kontrollierten Nutzung sichtbar macht.

\section{Abschließendes Fazit}
\label{sec:conclusion_final}

Die Arbeit zeigt, dass die Evaluation KI-gestützter Softwareentwicklung im Vibe-Coding-Kontext eine klare Trennung zwischen Intention, Implementationsartefakt und testgebundener Bestätigung erfordert.
Die Kombination aus Pipeline-Logik, Artefaktanalyse und Chat-Auswertung erlaubt eine konsistente, methodisch abgesicherte Zusammenführung der Befunde (Kapitel~6) und ihrer Einordnung (Kapitel~7).

Im Ergebnis ist der Nutzen im untersuchten Brownfield-Setting nicht als automatische Eigenschaft eines Tools zu lesen, sondern als Zusammenspiel von Integrationsaufwand, Validierungsarbeit und Nutzungstypen im Kompetenzkontext.

\noindent\textbf{Absicherung (Group comparability \& Confounder).} Die Schlussfolgerungen sind an die in Kapitel~6 berichteten Proxys (Pipeline-Status, Code-/Test-Surrogate, Chat-Muster) gebunden; mögliche Konfundierungen und die konservative Lesart werden in Abschnitt~\ref{sec:discussion_scope} (Kapitel~7) begründet. Entsprechend begründen die Befunde keine kausalen Effekte und keine generalisierbaren Rangfolgen zwischen Kompetenzgruppen.
Die Aussagen beziehen sich auf Branches als Analyseeinheiten und auf aggregierte Muster, nicht auf eine Leistungsbewertung einzelner Personen.

\noindent\textbf{Take-Home-Message.} Für KI-gestützte Softwareentwicklung sind singuläre Produktivitäts- oder Proxy-Metriken nur begrenzt belastbar, weil Integrations- und Validierungsarbeit zentral bleibt (Kapitel~6/7). Robust wird die Evaluation durch die konsequente Trennung und Triangulation von Intention (Attempted), Implementationsartefakt (Implemented\textsubscript{static}) und testgebundener Bestätigung (Verified). Kompetenz wird dabei durchgehend als Analyseperspektive genutzt, um unterschiedliche Nutzungs- und Validierungsmodi einzuordnen\,-- ohne Leistungsrangfolgen abzuleiten. Die Kernaussage ist: Belastbare Schlussfolgerungen entstehen erst aus der gemeinsamen Betrachtung dieser drei Ebenen.


\printbibliography

\end{document}
